\documentclass[10pt,a4paper]{article}

% ----------------- Paquetes -------------------%

\usepackage[T1]{fontenc}
\usepackage[utf8]{inputenc}
\usepackage[a4paper,margin=2.5cm]{geometry}
\usepackage{mathptmx}             
\usepackage{changepage} 
\usepackage{setspace}
\usepackage[spanish]{babel}
\usepackage[labelfont=bf,labelsep=space]{caption}
\usepackage{amssymb}
\usepackage{amsmath}
\usepackage{ebgaramond}
\usepackage{placeins}
\usepackage{etoolbox}
\usepackage{setspace}
\usepackage{parskip} 
\usepackage{tcolorbox}
\usepackage{needspace}
\usepackage{xparse}
\usepackage{ocgx2}
\usepackage{hyperref}
\usepackage{hypcap}


%%%%%%%%%%%%%%%%%%%%%%%%%%%%%%%%%%%%%%%%%%%%%%%%%%%%%%%%%%%%%%%%
% --- Fija primero tus valores globales "buenos" ---
\setstretch{1.2}                 % Interlineado global
\setlength{\parskip}{0.7\baselineskip}
\setlength{\parindent}{0pt}
\newlength{\GlobalParskip}
\setlength{\GlobalParskip}{\parskip}

\newcounter{eqnum}[subsection]
\renewcommand{\theeqnum}{\arabic{eqnum}}
\newcounter{alphaitem}
\newcounter{numitem}

%% ------------------- Recuadros----------------------%%
\newcommand{\puntosclave}[1]{%
  \begin{tcolorbox}[
    colframe=black,
    title={Puntos clave},
    before upper={\setlength{\parindent}{0.5cm}\setlength{\parskip}{0.5\baselineskip}}
  ]
  #1
  \end{tcolorbox}
}

\newcommand{\modificaciones}[1]{
  \begin{tcolorbox}[
    colback=white,
    colframe=red,
    title={Modificaciones a realizar},
    before upper={\setlength{\parindent}{0.5cm}\setlength{\parskip}{0.5\baselineskip}}
  ]
  #1
  \end{tcolorbox}
}


%% ------------------- Preguntas TEST----------------------%%
% Opciones: radio (single-choice)
\NewDocumentCommand{\OptRadio}{m m m}{%
  \noindent\CheckBox[radio,name=#1,value=#2,width=1.2ex,height=1.2ex]{}%
  \;\textbf{#2.}~#3\par
}

% Opciones: checkbox (multi-choice)
\NewDocumentCommand{\OptCheck}{m m}{%
  \noindent\CheckBox[name=#1,width=1.2ex,height=1.2ex,bordercolor={0 0 0}]{}%
  \;#2\par
}

% Pregunta single-choice (radio)
% Args: 1=id, 2=enunciado, 3=opciones, 4=solución+justificación, 5=imagen (o {}), 6=origen
\NewDocumentCommand{\PreguntaTestSingle}{m m m m m m}{%
  \par\medskip
  \noindent\textbf{#1.}~#2\par\smallskip

    % Imagen (si no es {})
    \IfBlankTF{#5}{}{%
      \begin{center}
        \includegraphics[width=0.75\linewidth]{#5}
      \end{center}
    }

    \begin{Form}
      #3
    \end{Form}

    \par\smallskip
    \switchocg{sol-#1}{\fbox{Mostrar / ocultar solución}}%
    \hfill{\footnotesize #6}\par\smallskip

    \begin{ocg}{Solución}{sol-#1}{0}
      \noindent #4
    \end{ocg}
  \par\medskip
}

% Pregunta multi-choice (checkboxes)
\NewDocumentCommand{\PreguntaTestMulti}{m m m m m m}{%
  \par\medskip
  \noindent\textbf{#1.}~#2\par\smallskip

    % Imagen (si no es {})
    \IfBlankTF{#5}{}{%
      \begin{center}
        \includegraphics[width=0.75\linewidth]{#5}
      \end{center}
    }
    \begin{Form}
      #3
    \end{Form}

    \par\smallskip
    \switchocg{sol-#1}{\fbox{Mostrar / ocultar solución}}%
    \hfill{\footnotesize #6}\par\smallskip

    \begin{ocg}{Solución}{sol-#1}{0}
      \noindent #4
    \end{ocg}
  \par\medskip
}


%% ------------------- Listas----------------------%%
    % --- Lista alfabética ---
    \newenvironment{la}
      {\begin{list}{\alph{alphaitem})}{%
          \usecounter{alphaitem}%
          \setlength{\leftmargin}{1cm}%
          \setlength{\parskip}{3pt}% 
          \setlength{\itemsep}{0pt}% ← espacio entre ítems
          \setlength{\parsep}{0pt}%  ← párrafos dentro de un ítem
      }}
      {\end{list}}
      
    % --- Lista numérica ---
    \newenvironment{lnum}
      {\begin{list}{\arabic{numitem}.}{%
          \usecounter{numitem}%
          \setlength{\leftmargin}{1cm}%
          \setlength{\parskip}{3pt}% 
          \setlength{\itemsep}{0pt}% ← espacio entre ítems
          \setlength{\parsep}{0pt}%  ← párrafos dentro de un ítem
      }}
      {\end{list}}
      
% ------------------- Imágenes----------------------%
    \graphicspath{{Img/}}% Rutas por defecto 
    % --- Imagen adaptativa ---
    % Registros de dimensión definidos una sola vez
    \newdimen\imgcap
    \newdimen\imgmin
    \newdimen\imgmax
    \newdimen\imgremain
    \newdimen\imgh
    \newdimen\imgneed
    
    \newcommand{\imagenfit}[3]{%
      \addvspace{10pt}% separación antes
      \par\begingroup
        % Parámetros de composición (asignaciones, no \newdimen)
        \imgcap=1\baselineskip            % alto estimado de título+fuente
        \imgmin=0.15\textheight
        \imgmax=0.15\textheight
        \imgremain=\pagegoal
        \ifdim\pagegoal=\maxdimen \imgremain=0pt\fi
        \advance\imgremain by -\pagetotal
        \advance\imgremain by -\imgcap
        % Decisión de altura destino
        \ifdim\imgremain<\imgmin
          \newpage
          \imgh=0.20\textheight
        \else
          \imgh=\imgremain
          \ifdim\imgh>\imgmax \imgh=\imgmax \fi
        \fi
        % Reservar espacio para evitar cortes del bloque completo
        \imgneed=\imgh \advance\imgneed by \imgcap
        \Needspace{\imgneed}
        % Cuerpo
        \noindent\begin{minipage}{\textwidth}
          \centering
          \captionof{figure}{\textemdash\ #2}\par\vspace{0.3em}% título arriba
          \includegraphics[width=\linewidth,height=\imgh,keepaspectratio]{\detokenize{#1}}\par
          \vspace{0.3em}\small\textit{Fuente: }#3% fuente abajo
        \end{minipage}%
      \endgroup
      \par\addvspace{10pt}%
    }

%------------ Bloque de RECUADRO ESPECIAL -------------%
    \newenvironment{specialblock}[1]{%
      \begin{quote} % crea sangría a izquierda y derecha
      \small % texto más pequeño
      \noindent\textbf{#1}\\[-0.3em] % título en negrita
      \rule{\linewidth}{0.4pt}\\[0.6em]}{
      \end{quote}}
      
%-------------------------- Portada -----------------------%
\newcommand{\oldtitle}[1]{\def\OldTitle{#1}}
\newcommand{\changerationale}[1]{\def\ChangeWhy{#1}}

\newcommand{\changebox}{%
\begin{tcolorbox}[title={Historial de cambios del tema}]
\textbf{Título anterior:} \OldTitle\par
\textbf{Motivo del cambio:} \ChangeWhy
\end{tcolorbox}
}

%-------------------------- Author Font -----------------------%
    \newcommand{\authorfont}[1]{{\fontfamily{EBGaramond-TLF}\selectfont #1}}

%-------------------------- Anexos -----------------------%
    \makeatletter
    \newcounter{anexo}
    \renewcommand{\theanexo}{\Roman{anexo}}
    \newcommand{\anexo}[1]{%
      \clearpage
      \phantomsection
      \refstepcounter{anexo}%
      \noindent\textbf{Anexo \theanexo.- }#1\par\medskip
      \addcontentsline{stc}{section}{Anexo \theanexo.- #1}%
    }
    \makeatother

    %-------------------------- Lista de figuras (personal) -----------------------%
\makeatletter
\newcommand{\MyListOfFigures}{%
  \clearpage
  \phantomsection
  {\centering\bfseries\Large Índice de figuras\par}%
  \medskip
  % Entrada en nuestro índice de contenido (stc)
  \addcontentsline{stc}{section}{Índice de figuras}%
  % Recuperación del listado (sin cabecera propia)
  \begingroup
    \parindent=0pt\parskip=2pt
    \@starttoc{lof}%
  \endgroup
}
\makeatother

%-------------------------- Bibliografía -----------------------%
\makeatletter
% Contador para las referencias
\newcounter{myref}
% Cabecera de la bibliografía, entra en el índice 'stc'
\newcommand{\MyBibliography}{%
  \clearpage
  \phantomsection
  {\centering\bfseries\Large Bibliografía y referencias\par}%
  \medskip
  \addcontentsline{stc}{section}{Bibliografía y referencias}%
  \setcounter{myref}{0}%
}
\newcommand{\MyBibItem}[4]{%
  \refstepcounter{myref}%
  \noindent[\themyref]\ #1.\ \emph{#2}. #3. #4\par
}
\makeatother

%--------- Establecer cuatro niveles de índice --------%
    \setcounter{tocdepth}{4}   
    \setcounter{secnumdepth}{4}% numera hasta \paragraph

%%%%%%%%%%%%%%%%%%%%%%%%%%%%%%%%

\makeatletter

    %--------- Interlineado Global --------%
    \edef\GlobalBaseLineStretch{\baselinestretch}
    
    %--------- Espaciado de seccinones --------%
    \renewcommand\section{\@startsection{section}
      {1}{\z@}
      {2.5ex \@plus 1ex \@minus .2ex} % espacio antes
      {1.5ex \@plus .2ex}             % espacio después
      {\normalfont\Large\bfseries}    % formato del título
    }
    \renewcommand\subsection{\@startsection{subsection}{2}{\z@}
      {3ex \@plus 0.8ex \@minus .2ex}
      {1.5ex \@plus .2ex}
      {\normalfont\large\bfseries}
    }
    \renewcommand\subsubsection{\@startsection{subsubsection}{3}{\z@}
      {3ex \@plus 0.5ex \@minus .2ex}
      {1ex \@plus .2ex}
      {\normalfont\normalsize\bfseries}
    }
    \renewcommand\paragraph{\@startsection{paragraph}{4}{\z@}%
    {-2ex \@plus -0.5ex \@minus -.2ex}% espacio antes
    {0.8ex \@plus .2ex}% espacio después
    {\normalfont\normalsize\itshape}}% ← cursiva en lugar de negrita

    %--------- Ecuaciones --------%   
    % Variables globales para almacenar títulos y notas
    \gdef\leq@current@section{}%
    \gdef\leq@current@subsection{}%
    \gdef\leq@last@printed@section{}%
    \gdef\leq@last@printed@subsection{}%
    
    % Comando para añadir nota (invisible en el cuerpo)
    \newcommand{\eqnote}[1]{%
      \addcontentsline{leq}{leqnote}{#1}%
    }
    
    % Comando para ecuaciones
    \newcommand{\eqblock}[2]{%
      % Verificar si necesitamos imprimir el título de sección
      \ifx\leq@current@section\leq@last@printed@section\else
        \addcontentsline{leq}{leqsection}{\leq@current@section}%
        \global\let\leq@last@printed@section\leq@current@section
        \gdef\leq@last@printed@subsection{}% Reset subsección al cambiar sección
      \fi
      % Verificar si necesitamos imprimir el título de subsección
      \ifx\leq@current@subsection\leq@last@printed@subsection\else
        \ifx\leq@current@subsection\@empty\else
          \addcontentsline{leq}{leqsubsection}{\leq@current@subsection}%
          \global\let\leq@last@printed@subsection\leq@current@subsection
        \fi
      \fi
      % Ecuación
      \refstepcounter{eqnum}%
      \begin{equation*}
        #1\tag{E.\theeqnum}
      \end{equation*}%
      \addcontentsline{leq}{leqt}{[E.\theeqnum] -- #2}%
      \addcontentsline{leq}{leqm}{#1}%
    }
    
    %%% ----- Índice de ecuaciones ----------%%%
    % Tamaño para []
    \newcommand{\leqIndexFont}{\normalsize}
    \newcommand{\leqTitleFont}{\tiny}
    \newcommand{\leqNoteFont}{\tiny}
    % Cabecera de sección 
    \newcommand*\l@leqsection[2]{%
      \addvspace{25pt}% Mayor separación antes de nueva sección
      \noindent\leqIndexFont
      \makebox[\linewidth]{\rule{.38\linewidth}{0.4pt}\ \ \textbf{  \MakeUppercase{#1}}\ \ \rule{.38\linewidth}{0.4pt}}%
      \par\vspace{-10pt}
    }
    % Cabecera de subsección 
    \newcommand*\l@leqsubsection[2]{%
      \par\addvspace{15pt}%
      \noindent\leqIndexFont #1
      \par\vspace{-7pt}
    }
    % Título de ecuación 
    \newcommand*\l@leqt[2]{%
    \par\addvspace{10pt}%
      \par\noindent\leqTitleFont #1\par
    }
    % Miniatura de ecuación
    \newcommand*\l@leqm[2]{%
      \vspace{0.3\baselineskip}%
      \par\scriptsize\noindent\hspace*{1.5em}\(#1\)\par%
    }
    % Nota de ecuación 
    \newcommand*\l@leqnote[2]{%
      \vspace{0.2\baselineskip}%
      \par\noindent\leqNoteFont\hspace*{1.5em}\textbf{Nota.--} #1\par\vspace{0.5\baselineskip}%
    }
    % Generación del índice
    \newcommand{\mylistofequations}{%
      \clearpage
      \phantomsection
      % Título visual de la lista (ajústalo a tu gusto):
      {\centering\bfseries\Large Índice de ecuaciones\par}
      % Entrada en nuestro índice 'stc':
      \addcontentsline{stc}{section}{Índice de ecuaciones}%
      % Llamada a tu lista real (usa el nombre exacto de tu macro):
      \@starttoc{leq}%
    }
    
    %--------- Captura de títulos de sección --------%
    \let\leq@original@section\section
    \let\leq@original@subsection\subsection
    % Flag para saber si estamos en una sección asterisco
    \newif\ifLeqStarSection
    \LeqStarSectionfalse
    
    % Redefinir \section CON CONTROL DE ASTERISCO
    \renewcommand{\section}{%
      \@ifstar{\leq@section@star}{\@dblarg\leq@section@nostar}%
    }
    \def\leq@section@star#1{%
      \global\LeqStarSectiontrue
      \gdef\leq@current@section{#1}%
      \gdef\leq@current@subsection{}%
      \leq@original@section*{#1}%
      \global\LeqStarSectionfalse
    }
    \def\leq@section@nostar[#1]#2{%
      \global\LeqStarSectionfalse
      \gdef\leq@current@section{#2}%
      \gdef\leq@current@subsection{}%
      \leq@original@section[#1]{#2}%
    }
    % Redefinir \subsection
    \renewcommand{\subsection}{%
      \@ifstar{\leq@subsection@star}{\@dblarg\leq@subsection@nostar}%
    }
    \def\leq@subsection@star#1{%
      \global\LeqStarSectiontrue
      \gdef\leq@current@subsection{#1}%
      \leq@original@subsection*{#1}%
      \global\LeqStarSectionfalse
    }
    \def\leq@subsection@nostar[#1]#2{%
      \global\LeqStarSectionfalse
      \gdef\leq@current@subsection{#2}%
      \leq@original@subsection[#1]{#2}%
    }

    % ----- Restauración de valores Globales de estilo ----------%
    % Flags
    \newif\ifInFootnote
    \InFootnotefalse
    \pretocmd{\@footnotetext}{\InFootnotetrue}{}{}
    \apptocmd{\@footnotetext}{\InFootnotefalse}{}{}
    % Profundidad de listas (soporta anidadas)
    \newcount\MyListDepth \MyListDepth=0
    \AtBeginEnvironment{la}{\advance\MyListDepth by 1}
    \AfterEndEnvironment{la}{\advance\MyListDepth by -1}
    \AtBeginEnvironment{lnum}{\advance\MyListDepth by 1}
    \AfterEndEnvironment{lnum}{\advance\MyListDepth by -1}
    \AtBeginEnvironment{itemize}{\advance\MyListDepth by 1}
    \AfterEndEnvironment{itemize}{\advance\MyListDepth by -1}
    \AtBeginEnvironment{enumerate}{\advance\MyListDepth by 1}
    \AfterEndEnvironment{enumerate}{\advance\MyListDepth by -1}
    % Restauración diferida: hazlo al INICIO del próximo párrafo real
    \newif\if@pendingRestore
    \@pendingRestorefalse
    \newcommand{\RequestRestore}{\global\@pendingRestoretrue}
    \everypar{%
      \if@pendingRestore
        \global\@pendingRestorefalse
        \ifInFootnote\else
          \ifnum\MyListDepth=0
            \setlength{\parskip}{\GlobalParskip}%
            \setstretch{\GlobalBaseLineStretch}%
          \fi
        \fi
      \fi
    }
    % Al final de bloques:
    \AfterEndEnvironment{equation}{\RequestRestore}
    \AfterEndEnvironment{equation*}{\RequestRestore}
    \AfterEndEnvironment{la}{\RequestRestore}
    \AfterEndEnvironment{lnum}{\RequestRestore}
    % Al final de macros:
    \apptocmd{\eqblock}{\RequestRestore}{}{}
    \apptocmd{\imagenfit}{\RequestRestore}{}{}
    
\makeatother

% ===================== ToC propio con 4 niveles (único bloque) =====================
\makeatletter

% --- Escritores: añadimos una línea al .stc en cada cabecera numerada ---
% Guardamos las versiones actuales para llamar al comportamiento normal
\let\MyTOC@orig@section\section
\let\MyTOC@orig@subsection\subsection
\let\MyTOC@orig@subsubsection\subsubsection
\let\MyTOC@orig@paragraph\paragraph

% SECTION
\renewcommand{\section}{%
  \@ifstar{\MyTOC@section@star}{\@dblarg\MyTOC@section@nostar}%
}
\def\MyTOC@section@star#1{%
  \MyTOC@orig@section*{#1}% sin entrada al .stc
}
\def\MyTOC@section@nostar[#1]#2{%
  \MyTOC@orig@section[{#1}]{#2}% comportamiento normal (escribe al .toc)
  % Escribimos también nuestra línea al .stc (texto con \numberline)
  \addcontentsline{stc}{section}{\protect\numberline{\thesection}#2}%
}

% SUBSECTION
\renewcommand{\subsection}{%
  \@ifstar{\MyTOC@subsection@star}{\@dblarg\MyTOC@subsection@nostar}%
}
\def\MyTOC@subsection@star#1{%
  \MyTOC@orig@subsection*{#1}%
}
\def\MyTOC@subsection@nostar[#1]#2{%
  \MyTOC@orig@subsection[{#1}]{#2}%
  \addcontentsline{stc}{subsection}{\protect\numberline{\thesubsection}#2}%
}

% SUBSUBSECTION
\renewcommand{\subsubsection}{%
  \@ifstar{\MyTOC@subsubsection@star}{\@dblarg\MyTOC@subsubsection@nostar}%
}
\def\MyTOC@subsubsection@star#1{%
  \MyTOC@orig@subsubsection*{#1}%
}
\def\MyTOC@subsubsection@nostar[#1]#2{%
  \MyTOC@orig@subsubsection[{#1}]{#2}%
  \addcontentsline{stc}{subsubsection}{\protect\numberline{\thesubsubsection}#2}%
}

% PARAGRAPH
\renewcommand{\paragraph}{%
  \@ifstar{\MyTOC@paragraph@star}{\@dblarg\MyTOC@paragraph@nostar}%
}
\def\MyTOC@paragraph@star#1{%
  \MyTOC@orig@paragraph*{#1}%
}
\def\MyTOC@paragraph@nostar[#1]#2{%
  \MyTOC@orig@paragraph[{#1}]{#2}%
  % Si no numeras los \paragraph, cambia \theparagraph por texto plano sin \numberline
  \addcontentsline{stc}{paragraph}{\protect\numberline{\theparagraph}#2}%
}

% --- Impresor del índice propio (lee el .stc) ---
\newcommand{\MyTOC}{%
  \par\bigskip
  {\centering\bfseries\Large Índice de contenido\par}%
  \medskip
  \begingroup
    \parindent=0pt\parskip=2pt
    \@starttoc{stc}%
  \endgroup
}

\makeatother
% ===================== fin ToC propio =====================


%%%%%%%%%%%%%%


% --- Datos del tema ---
\title{Tema 3.A.1 -- Objeto y método de la ciencia económica\\
\large Cuestiones y debates actuales, con especial referencia a la economía conductual}
\author{Marcelo Ortega}
\date{Noviembre 2025}
\newcommand{\TemaCode}{3A01}

\oldtitle{Tema 3.A.1. Objeto y método de la ciencia económica}
\changerationale{Se explicita que el opositor no debería limitarse a realizar una revisión histórica y neutral, sino que debería incluir las presentes problemáticas y las controversias que están guiando el futuro de la Economía (como el Behavioural Economics).}


\begin{document}


\maketitle
\changebox

\newpage

% --- Página de resumen y modificaciones ---
\puntosclave{ %% Escribir puntos clave Apuntes CECO
    Si los temas en la oposición son siempre abiertos y es labor del opositor realizar las labores de investigación y profundización ( con las fuentes adecuadas) para alcanzar versiones auténticamente propias de los mismos, en este caso, lo es si cabe más aún.\footnote{Cuestiones que el autor ha preferido no introducir en el tema. pero que pueden resultar de interés son. sin ánimo de exhaustividad: (i) el concepto de homo economicus y el Rational Choice (qué es, porqué su recurso es extendido en Economía pese a las criticas generales desde otras ciencias sociales y desde la propia disciplina); (ii) los principales debates entre escuelas de pensamiento económico (en especial, la existencia y optimalidad de equilibrios descentralizados, y por tanto el grado deseable de intervención pública en la economía); (iii) las virtudes y limitaciones estadísticas de la econometría; (iv) la proliferación de herramientas de análisis de datos y su impacto en la acotación de la ciencia económica, frente a otras como la sociología, la Historia o la psicología; o, (v) la relación de la Economía con la Filosofia y la Ética.}
    }
\vspace{1cm}

\modificaciones{
    Añadir: Muy importante \href{https://www.exploring-economics.org/en/orientation/#discover}{\textit{Compare the perspectives of economics}. Exploring Economics}

    Añadir: \textbf{¿Es la economía una Ciencia?} -> Estructuras la respuesta invertida: ¿Qué es una \textit{ciencia}? ¿Qué importancia real tiene el comportamiento holístico en el campo científico? ¿Cómo se estructura la investigación científica? ¿Por qué no es una ciencia? La investigación científica se sustenta sobre la completitud de sus sitemas lo que permite a cada investigador dar un paso adelante y avanzar la frontera del conocimiento. Nosotros, los economistas, nos seguimos peleando sobre el rol que ocupa el dinero en nuestra economía, ¿cómo podemos atribuirle carácter científico a un campo que sigue enzarzado en discusiones elementales? En economía encontrará ricardiamos, neo paretianos, keynesianos, monetaristas, neokeynesianos, defensores de la síntesis, etc. En una ciencia encontrará científicos especializados en una de las areas del conocimiento objeto de estudio. La investigación económica se sustenta sobre el sistema de creencias que tiene cada investigador. Eso no quiere decir que no tenga validez, sino que le atribuimos validez de una manera erronea: le damos validez a nuestras creencias y no a nuestra investigación.

}

\newpage
\MyTOC
\newpage

\mylistofequations
\clearpage

\MyListOfFigures
\clearpage


\pagenumbering{arabic}
\setcounter{page}{1}


\section{Introducción}\label{intro}
Entender de donde venimos resulta fundamental para saber hacia donde vamos. 

Aunque pueda parecer \textit{naive} es de una importancia fundamental y, más aún, cuando nos movemos en el mundo de la epistemología: la Teoría del Conocimiento. 

Un ejemplo parádigmatico de esto resulta el Premio Nobel de Economía de 2024, el concepto de \textit{destrucción creativa}. Un autor puede encuadrarse en una época, resultar ser el contorno entre épocas o precisamente las ideas germinales de cambio que motivan el nacimiento de otra época, pero lo que es innegable es que ese autor es fruto del pasado. 

Si consideramos el conocimiento como una escalera autoncontenida y en desarrolla, la \textit{destrucción creativa} nos permite comprender el valor de dejar que las ideas florezcan y destruyan a sus antecesoras; pero ese escalón desde el cual brotan resulta inalcanzable sin la presencia de todos aquellos que aportaron su pieza a ésta. 

\imagenfit{DestruccionCreativa.png}{Ilustración: Destrucción creativa}{\href{https://www.nobelprize.org/prizes/economic-sciences/2025/press-release/}{Nobel Prize Outreach 2025. Nov 2025.}}

\imagenfit{PrescriptiveKnowledgeVSPropositionalKnowledge.png}{Ilustración: Conocimiento Prescriptivo y Conocimiento Proposicional}{\href{https://www.nobelprize.org/prizes/economic-sciences/2025/press-release/}{Nobel Prize Outreach 2025. Nov 2025.}}


\subsection{Contextualización}\label{context}

\subsubsection{Relevancia}\label{importance}


\subsubsection{Historia}\label{history}
En nuestros días, no existe una definición inequívoca ni canónica de economía, pero muchos de los significados que denota hoy el término economía están conectados con el sentido en el que JENOFONTE\footnote{En el 400 a. C, un discípulo de SÓCRATES llamado JENOFONTE escribió un diálogo socrático conocido como \textit{Oikonomikos}. JENOFONTE utilizaba el término \textit{oikonomia} para denotar la buena gestión de la casa y la hacienda por medio de la buena llevanza de los asuntos diarios y la tenencia ordenada de las cuentas de los negocios familiares.} utilizó el término. Así, economía se aplica hoy a la gestión de la casa y la hacienda personal, pero también de una sociedad mercantil, de una institución estatal o de cualquier otro tipo de organizaciones. También, el término economía se utiliza como sinónimo de buena gestión de un recurso escaso, de ahorro, y de adopción de decisiones que nos acercan a un resultado considerado óptimo. En el significado que resulta más directamente relevante a este tema de esta oposición, la palabra economía hace referencia a un corpus de conocimiento que tiene por objetivó entender y predecir fenómenos humanos diversos como las decisiones de compra y venta, la relación entre el paro e inflación, o la relación entre los directivos, los trabajadores y los accionista de una empresa, entre muchos otros fenómenos.

\subsubsection{Problemática}\label{problem}
Las variadas concepciones de la economía se concretan en respuestas a dos preguntas esenciales: 
\begin{lnum}
    \item ¿Qué estudia la economía?; y,
    \item ¿Cómo lo estudia?
\end{lnum}

O de otro modo: ¿qué objetivo tiene el estudio de la economía? ¿qué trata de entender y predecir?\footnote{Y habiendo contestado aquellas, responder a: ¿cómo avanzar en ese conocimiento? ¿cómo sistematizar el conocimiento adquirido sobre los fenómenos económicos a estudiar?, ¿de qué manera organizarlo y transmitirlo? ¿de qué forma optimizar la obtención de nuevas explicaciones de fenómenos observados?, ¿cómo mejorar la predicción de fenómenos económicos futuros?, ¿qué métodos de obtención de conocimiento económico emergenen los últimos tiempos?, ¿qué relaciones existen entre ellos?, ¿cómo valorar unos frente a otros?}

\subsection{Estructura de la exposición}\label{structure}
Estructuraremos nuestra exposición en cuatro partes:
\begin{lnum}
    \item Objeto de la ciencia económica. Es decir, ¿qué estudia la economía?
    \item Método de la ciencia económica. Analizaremos las diferentes formas de estudiar el objeto de la economía, atendiendo a los aspectos epistemológicos generales de la economía, los principales debates metodológicos entre economistas, así corno las cualidades deseables de los modelso económicos.
    \item Fronteras de investigación. Entre otros, las aportaciones metodológicas de la economía conductual (behavioral economics).
\end{lnum}










\clearpage
\section{Ciencia Ecónomia: Objeto}
\puntosclave{
    El objeto de la Economía ha variado a lo largo de la Historia. Los economistas han centrado sus esfuerzos en explicar y valorar diferentes tipos de fenómenos, tales como el buen gobierno desde el punto de vista de quién está al frente de una organización estatal, o la gestión de los asuntos diarios.

    Ha existido una tensión de largo plazo entre dos economía como una ciencia que valore lo que debe ser, y una ciencia que explique qué es.

    Las herramientas metodológicas de la economía se han aplicado al estudio de fenómenos que inicialmente se consideraban objeto de otras disciplinas sociales.
}

¿Qué tipo de fenómenos estudia la economía? ¿Cuál es el objetivo de los economistas cuando tratan de entender y predecir fenómenos económicos? Las respuestas a estas preguntas han sido diversas y muchas de ellas, si no todas, siguen siendo relevantes a los fenómenos económicos actuales.\footnote{El análisis de esta cuestión desde una perspectiva histórica permite poner en contexto la evolución de los objetos de la economía con otras disciplinas, así como con el contexto social y político, [\authorfont{Ver Anexo \ref{anx:historiaobjeto}}]}

\subsection{Teoría de la Elección: la Decisión}
La Ciencia Económica es, esencialmente, una Teoría de la Elección. La elección constituye el problema económico primigenio, el pecado original. 

¿Quiero comprarme esa bolsa de pipas? ¿Debo comprar un coche ahora o esperarme y permitirme uno mejor? ¿Cuántos años debo dedicar al estudio? ¿Es una buea idea empezar las oposiciones? ¿Debo comprarme una casa o es mejor alquilar? ¿Cuántos hijos tener? ¿Cuánto dinero debo ahorrar para el futuro? ¿Debería casarme en régimen de gananciales o en separación de bienes? ¿En cuánto debemos fijar la prestación por desempleo? ¿Cómo debemos luchar contra la pobreza infantil? ¿Cuánto dinero debemos asignar a la renovación de X tramo de autopista? ¿Debemos construirla en primer lugar?

Todas estas son preguntas que a diario enfrenta uno, muchos o todo el conjunto de personas que conforman nuestro mundo. Pero, para lo que nos interesa, lo que sí sabemos es que el objeto de la ciencia económica es la elección/decisión. No obstante, el enfoque que se adopte respecto a este objeto es completamente discrecional y en función de la vertiente económica a la que nos afiliemos. Lo que queremos decir es que no se puede especificar con generalidad cómo enfoca el estudio del objeto: ¿Cómo toma esta decisión? ¿Cuál debería tomar? ¿Por qué toma esa decisión? Son enfoques distintos y todos ellos conforman la ciencia económica.



\textcolor{red}{\textbf{OJO!} Revisar a partir de aquí y reconciliar con la nueva estructura (ha cambiado mucho).}

La teoría económica moderna examina los mecanismos que utilizan las sociedades contemporáneas para resolver los problemas que se derivan de la escasez relativa. Centra la atención principalmente en los procesos del mercado, que han ocupado el lugar de la Iglesia, la tradición y el Estado como principal mecanismo de asignación de los recursos.\footnote{Sin embargo, estos mecanismos no son mutuamente excluyentes, como tampoco ha sido lineal la transición de las economías basadas en la tradición, el Estado y la Iglesia a una economía de mercado; no todas las sociedades del mundo han participado. En algunas zonas, que abarcan casi todo un continente, la actividad económica sigue estando dominada por el pasado. En algunas sociedades, se ha pasado de una economía feudal anterior al mercado a una economía autoritaria moderna, en la que el Estado asigna los recursos. Por ejemplo, a principios de la década de 1900, algunas sociedades adoptaron un sistema de planificación central, que implicaba el control estatal de la asignación de los recursos. En Europa Oriental, se observan movimientos para pasar de una economía autoritaria a una economía de mercado, cuyos resultados son inciertos.

Decir que el mercado es el principal mecanismo de asignación no es decir que sea el único. Las sociedades modernas de mercado utilizan la fuerza, la tradición y la autoridad, además de los mercados. En Europa y en Norteamérica, las fuerzas sociales y políticas influyen continuamente en la asignación basada en el mercado.

A modo de ejemplo, es innegable que los mercados tuvieron una gran relevancia en las ciudades del imperio romano. Pero de ello no puede deducirse que la economía de Roma fuese un sistema de mercado, porque éste no era el principal mecanismo de asignación de los recursos. En cualquier caso, aquellos mercados no actuaban como lo harían en las economías capitalistas, pues el comercio a largas distancias del mundo romano no respondió a estímulos económicos (el principio ricardiano de la especialización regional según las ventajas comparativas), sino a decisiones políticas, tomadas en Roma, sobre la distribución regional del excedente captado por el Imperio. Más que una economía de mercado, la de Roma fue una mezcla de imperio tributario y de sistema esclavista. Esto significa que el comportamiento de los comerciantes, latifundistas y agricultores del mundo romano no puede evaluarse con los conceptos de la racionalidad económica del capitalismo. La economía del mundo romano atendía a su propia lógica, que estaba adaptada al contexto político, institucional y social de entonces.

La teoría económica moderna aún está tratando de comprender las interrelaciones entre las fuerzas económicas, sociales y políticas. Ha centrado la atención en el modo en que funcionan las fuerzas del mercado, concentrando sus esfuerzos en averiguar cómo asignan los mercados los recursos escasos y cuáles son las fuerzas que determinan el nivel de producción económica y su crecimiento. Pero el pensamiento económico va más allá de esas cuestiones. Como vemos, muchos de los interrogantes que suscitaron los primeros autores anteriores a la aparición del mercado se referían a cuestiones filosóficas y éticas más generales que ayudan a situar en perspectiva el pensamiento económico moderno.

Independientemente de cuál sea el mecanismo que utilice la sociedad para asignar los recursos, la cruda realidad de la escasez obliga a dejar algunos deseos sin satisfacer; por tanto, la raíz del problema de la escasez es la cuestión de la equidad y la justicia. Los mecanismos de asignación de los recursos deciden quién recibe los recursos y quién no.} En la mayoría de sociedades desarrolladas, la asignación de recursos se realiza a través del mercado, es decir, a través de las decisiones descentralizadas de millones de agentes. Por ello, normalmente, la economía se preocupa de cuestiones relacionadas con el funcionamiento del mercado, y esto lleva a preguntarse cómo toman decisiones los agentes, es decir, el estudio del comportamiento de los distintos agentes que forman el sistema económico. De esta forma, gira entorno a varias ideas fundamentales:
\begin{lnum}
    \item \textit{Recursos escasos}.\footnote{Otra posible crítica a la importancia de la escasez para definir qué es la economía es que se puede argumentar que todas las ciencias sociales parten en cierto modo de esa idea. Por ejemplo, si no hubiera limitación no existiría el derecho patrimonial, pues precisamente la delimitación de la propiedad y de las libertades individuales da origen a los problemas estudiados por esta rama.} De acuerdo con la tercera ley de GOSSEN (1854), la escasez de recursos es una condición previa al análisis económico. Es decir, existe un problema económico porque los recursos son escasos.\footnote{Posteriormente, LIPSEY, criticaría esta definición, afirmando que la economía no sólo consiste en la gestión de los recursos escasos, ya que considera que incluso cuando no son escasos pueden estar mal asignados. De este modo, según LIPSEY, la economía es el estudio del aprovechamiento óptimo de los recursos.}
    \item \textit{Fines alternativos}. Influenciados por la escasez, los agentes tienen que enfrentarse a disyuntivas y tomar decisiones, no sólo a nivel individual, sino también a nivel social (p.ej. disyuntiva entre eficiencia y equidad). La necesidad de elegir nos lleva a pensar en términos del coste de oportunidad (pensar en lo que renunciamos al tomar una decisión).
    \item \textit{Comportamiento de los individuos}. GARY BECKER, afirma que la economía es el estudio del comportamiento de los individuos que se caracterizan por tomar todo tipo de decisiones\footnote{BECKER recibió el Premio Nobel de Economía en 1992 «Por extender el dominio del análisis microeconómico hacia nuevos dominios del comportamiento y de las relaciones humanas, incluso más allá de los límites del mercado». El galardón sintetiza la forma de entender la forma de entender la disciplina por BECKER. La Economía es una forma de pensar que puede ser empleada como herramienta para explicar cualquier ámbito. De hecho, cuando recibe el premio Nobel, la charla lleva como título \textit{The Economic Way of Looking at Life}. 
    
    Esa concepción de la Economía como una forma de pensar susceptible de ser aplicada, lleva a BECKER a entrar a analizar con lógica económica algunos temas como el divorcio, el crimen, la adicción, la venta de órganos o la distribución de las tareas del hogar en la familia: se habla entonces de \textbf{imperialismo de la Economía}, pues estos campos no se consideraban tradicionalmente objeto de estudio por los economistas [\authorfont{Ver Anexo \ref{anx:historiaobjeto}}]. La aproximación del autor y su concepción de Economía como herramienta ha recibido fuertes críticas, pero también valoraciones favorables. 
    \begin{la}
        \item Entre los defensores del uso de los métodos de la ciencia económica para estudiar otros objetos distintos a los que inicialmente dieron lugar a la economía, es habitual afirmar que ninguna otra ciencia social ha logrado mejor capacidad predictiva, consistencia empírica, identificación explícita de los métodos aplicados y precisión formal que la economía, y por ello, no es indeseable aplicar las herramientas de la economía a fenómenos que se escapa de su objeto tradicional.
        \item Algunos de los argumentos contrarios al imperialismo de la economía apuntan a la ausencia de consideraciones éticas y morales en el análisis económico actual, la escasa fiabilidad de los modelos económicos para predecir a futuro, o la aplicación de herramientas metodológicas que no mejoran el conocimiento de la realidad, pero sí introducen una barrera de entrada con el objetivo de mantener las rentas que capturan los economistas en el diseño de políticas o en la academia.
    \end{la} 
    De hecho, la Economía de la Salud y la Economía de la Educación son actualmente dos ramas de la economía que nacen de los desarrollos de investigadores aplicando las herramientas propias de la Economía a estos ámbitos. 
    
    En definitiva, sí que podemos asegurar que la economía no sería hoy lo que es sin GARY BECKER. Un ejemplo de ello es el éxito del libro \textit{Freakonomics} (2005), del que STEVEN LEVITT (discípulo de BECKER) es coautor.
    } conscientemente intentando maximizar su bienestar y teniendo en cuenta los costes y beneficios de las distintas alternativas disponibles; y,
    \item \textit{Toma de decisiones consciente.} Ya sea desde una perspectiva racional, heurística o mediante cualquier técnica evaluativa, se considera el análisis de las decisiones individuales desde una perspectiva de sujetos activos.\footnote{Esta innovación, explicitada por GARY BECKER pero con larga tradición dentro de la economía resulta de gran importancia, pues permite realizar predicciones acerca de los efectos potenciales de distintas políticas públicas, por ejemplo sobre la tasa de fertilidad o el nivel de criminalidad en una sociedad. Otras disciplinas como el derecho interpretaban la actividad criminal como comportamientos irracionales o determinados por el entorno social y cultural, tratando por tanto a los sujetos como pasivos.}
\end{lnum}

Otros autores han tratado de describir la economía mediante otras lineas o han complementado la definición expuesto hasta ahora, sin embargo, el objeto de esta ciencia será el estudio de situaciones que compartan estas características (todas o alguna).\footnote{JOSÉ LUIS FERREIRA, en su libro Economía y Pseudociencia, define a los economistas como los que investigan en Economía y publican en revistas académicas sometiéndose a la revisión por pares (la revisión por pares o arbitraje (peer reviewing) es una evaluación usada para valorar trabajos escritos realizada por una o más personas con competencias similares a los productores del trabajo (expertos) pero que no forman parte del personal editorial del trabajo a evaluar, con el fin de asegurar la calidad, factibilidad y rigurosidad científica del trabajo. Funciona como una forma de autorregulación de miembros calificados de una profesión dentro del campo relevante. Los métodos de revisión por pares se utilizan para determinar los estándares de calidad técnica y científica, proporcionar credibilidad y corregir los artículos originales escritos por los investigadores), o los que sin investigar conocen y usan los resultados de las investigaciones hechas por los demás. Ello incluiría, por tanto, no sólo a académicos pero también a personas que están inmersos en la creación y coordinación de la política económica. La definición sí excluiría a personas cuya labor profesional es la dirección de empresas, la participación en los mercados financieros, la asesoría fiscal o la consultoría, pero no excluiría a personas que han utilizado la teoría económica en el campo de las finanzas, como por ejemplo, MARKOWITZ, MILLER o SHARPE todos ellos Premio Nobel de Economía en 1990. 

Por su parte, TIROLE (galardonado con el Premio Nobel de Economía en 2014 \textit{Por sus análisis sobre el poder y las regulaciones del mercado}) afirma que los economistas hacen algo análogo a los médicos: diagnosticar y proponer el mejor tratamiento posible (política económica) dado el conocimiento de la disciplina.}

Sin embargo, no sólo se analiza el funcionamiento del mercado, sino que también se estudian posibles intervenciones del Estado (economía pública y economía del bienestar) e incluso la asignación de recursos en economías intervenidas (p.ej. VON MISES, HAYEK y OSKAR LANGE\footnote{Podríamos proponer el planteamiento extremo suponiendo que el sector público interviene totalmente la economía (economía de planificación). Según LUDWIG VON MISES en una economía de planificación surgirían problemas de cálculo económico ya que en ausencia de un sistema de precios los distintos recursos de una sociedad no podrían ser asignados correctamente.
Frente a esto, OSKAR LANGE propone un socialismo de mercado, en el que la autoridad central anuncia precios y los individuos actúan de acuerdo a ellos. Mediante un mecanismo de prueba y error, aplicando la lógica de la optimalidad paretiana, una Oficina Central de Planificación ejercería de hecho las mismas funciones que el mercado. El argumento de LANGE convenció a muchísimos economistas y fue una fuerza clave detrás del renacimiento de la teoría del equilibrio general después de la Segunda Guerra Mundial. Autores como KENNETH ARROW querían tener un modelo matemático de cómo podría funcionar la economía de mercado precisamente porque querían construir un sistema económico alternativo.

HAYEK, en su trabajo \textit{The Use of Knowledge in Society} (1945), critica la posibilidad de socialismo de mercado. HAYEK explica que la mayoría de las discusiones sobre la eficiencia de un sistema económico parte de la hipótesis que conocemos las preferencias y la tecnologías. Sin embargo, dicha información no la posee la autoridad central, sino que la información de preferencias y tecnologías (lo que llama HAYEK el \textit{conocimiento}) está desperdigada entre todos los individuos. La mejor manera de hacer aflorar esta situación es mediante la interrelación descentralizada de los agentes en el mercado. Ello lleva a que los precios de mercado, al suministrar información sobre las preferencias y la escasez relativa de los bienes sea una buena manera según el austriaco de agregar la información.

Sin embargo, cuando existen problemas como externalidades, bienes públicos o información asimétrica, el sistema de precios puede agregar la información de manera deficiente. A pesar de ello, el trabajo de HAYEK tuvo, quizás de una manera menos visible pero no por ello menos cierta, una gran influencia en muchos economistas. Aunque la teoría moderna, altamente formal y matemática, se parezca poco a lo que HAYEK pensaba que era la mejor manera de teorizar, su herencia es indudable.

Por ejemplo, HURWICZ comenta que fueron las ideas de HAYEK las que le llevaron a pensar en temas como el diseño de mecanismos y la teoría de la compatibilidad de incentivos y que tratan precisamente de pensar cuándo y cómo obtener información dispersa de manera rigurosa para que la información descentralizada sea superior a la que llegaría en ausencia del mecanismo.

En definitiva, el fallo del sector público que subyace es la gran cantidad de información presente en los agentes de la que no dispone el sector público y le incapacita para actuar como planificador benevolente [\authorfont{Ver Tema 3.A.23}].}). 

\subsubsection{Enfoque: ¿positivo o normativo? ¿ideal o descriptivo?}
Qué debe hacer la ciencia económica:
\begin{la}
    \item ¿Debe limitarse a describir lo que es?; o,
    \item ¿Debe ir más allá, y tratar de determinar lo que debe ser?
\end{la}

Al analizar la evolución histórica de las distintas definiciones  del objeto de la economía [\authorfont{Ver Anexo \ref{anx:historiaobjeto}}], hallamos que el hilo conductor del debate acerca del objeto de la economía gira en torno al carácter positivo o normativo que ésta debe tomar. Siguiendo al profesor DAVID COLANDER:\footnote{Esta disyuntiva subyace en cierto grado a todas las definiciones históricas propuestas [\authorfont{Ver Aexo \ref{anx:historiaobjeto}}], y ha sido objeto de intensos debates:
\begin{la}
    \item El análisis moral del intercambio de los escolásticos medievales europeos trasluce un qué debe ser, y por ello, puede entenderse como una concepción normativista del objeto de la economía;
    \item En sentido similar, aunque de forma menos acusada, la concepción de la economía como ciencia del buen gobierno apunta en la dirección del normativismo;
    \item La ciencia económica como descubrimiento de leyes naturales tiende a lo contrario, y orienta el objeto de la economía hacia el qué es;
    \item El análisis marxista plantea una visión mixta ya que, si bien trata de aplicar un análisis sistemático  para descubrir una serie de leyes que caractericen el qué es, la finalidad transformadora de este análisis es explícita, y resulta en un diálogo constante entre el qué es y el qué debe ser;
    \item La concepción de MARSHALL, y posteriormente la de ROBBINS y SAMUELSON, proponen de nuevo un equilibrio: la economía describe lo que es cuando se está tratando de alcanzar un deber ser.
\end{la}
Es decir, tratan de describir cómo alcanzar un objetivo que se asume deseable -con diferentes grados de justificación-. Posteriormente, MILTON FRIEDMAN y otros autores hicieron notar la íntima relación que existe entre positividad y normatividad: ¿es acaso posible afirmar qué debe ser antes de conocer suficientemente el qué es? También en el siglo XX, la llamada Economía del Bienestar ha tenido por objeto aislar el qué debe ser mediante la identificación de los juicios de valor y las preferencias que hacen deseables unos estados de la sociedad respecto a otros.}
\begin{la}
    \item Si la economía es una ciencia positiva, debe tratar de reflejar la realidad de los fenómenos y dejar a un lado las consideraciones acerca de lo que debe ser. Dicho de otra forma, el único objetivo de la economía positiva es el conocimiento por conocimiento, dejando de lado en la medida de lo posible los juicios normativos;
    \item Si la economía es una ciencia normativa, su objetivo es dilucidar qué acciones deben tomar los sujetos analizados. La economía normativa es la rama filosófica de la economía e integra economía y ética. Por tanto, debe incorporar juicios de valor; y,
    \item Por otra parte, el arte de la economía\footnote{Término introducido por JOHN NEVILLE KEYNES, padre de JOHN MAYNARD KEYNES.} se ocupa de cuestiones relacionadas con la política económica. Relaciona economía positiva y normativa, y formula preguntas como la siguiente: \textit{Si éstos son mis objetivos y esta es la forma en que funciona la economía, ¿cuál es la mejor manera de lograr estos objetivos?}
\end{la}

\paragraph{Impacto sobre el método}
La distinción es importante porque la economía positiva y el arte de la economía tienen metodologías muy distintas.
\begin{la}
    \item La metodología de la economía positiva es formal y abstracta; trata de separar las fuerzas económicas de las fuerzas políticas y sociales; y,
    \item La metodología del arte de la economía es más compleja, ya que se ocupa de la política económica y debe abordar las relaciones entre la política, las fuerzas sociales y las fuerzas económicas. En el arte de la economía hay que integrar todas las dimensiones de un problema de las que se hace abstracción en la economía positiva.
\end{la}

\paragraph{¿Economía positiva o arte de la economía?}
¿Cuál debe ser el principal tema de interés de la economía? ¿La economía positiva o el arte de la economía?
Ésta es una cuestión que ha suscitado un interminable debate en la historia del pensamiento económico.
\begin{la}
    \item La escuela histórica alemana y la escuela marshalliana inglesa propugnan que se preste atención principalmente al arte de la economía. Para defender su postura se inspiran en la obra de ADAM SMITH;
    \item Los economistas ortodoxos modernos ponen el énfasis en la economía positiva y defienden su postura basándose en los escritos de DAVID RICARDO. En consonancia con esa idea, la mayoría de los escritos metodológicos modernos han girado en torno a la economía positiva.
\end{la}

\paragraph{¿Economía positiva o normativa?}
Este debate se conoce también como \textit{controversia FRIEDMAN-MYRDAL}):
\begin{la}
    \item En defensa de la economía normativas:
    \begin{lnum}
        \item GUNNAR MYRDAL\footnote{GUNNAR MYRDAL fue galardonado con el Premio Nobel de Economía en 1974 junto con FRIEDRICH HAYEK \textit{Por sus trabajos en la teoría del dinero y de las fluctuaciones y por su análisis de la independencia de los fenómenos económicos, sociales e institucionales}.
        } consideraba que los principales conceptos económicos se hallan cargados implícitamente de valor, por lo que defendía el predominio de la economía normativa.
        \item ATKINSON, \textit{(...) economics is essentially a moral science. Welfare economics should be a central part of the discipline.} Deben de incorporarse principios morales como la desigualdad, la libertad real de las personas para tomar decisiones, la igualdad de oportunidades, etc;
        \item KEYNES será también un economista moral/normativo. Su teoría tenía un claro fin: mejorar la vida de las personas, lo que pasaba por mejorar funcionamiento de las economías de mercado; y,
        \item KENNETH ARROW, por su parte, adopta una postura según la cual lo normativo puede dar lugar a lo positivo. ARROW buscaba diseñar un sistema económico alternativo más deseable que el de mercado (afirmación normativa) y, por ello, profundiza en las propiedades del equilibrio general, buscando un modelo matemático de cómo funciona la economía de mercado.
    \end{lnum}
    \item En defensa de la economía positiva:
    \begin{lnum}
        \item ROBBINS, "\textit{(...) Economics deals with facts; ethics with valuations}";
        \item  MILTON FRIEDMAN,\footnote{MILTON FRIEDMAN fue galardonado con el Premio Nobel de Economía en 1976 \textit{Por sus triunfos en el campo del análisis del consumo, la historia y teoría monetaria, y por su demostración acerca de la complejidad de la estabilización política}.} en cambio, defenderá en su ensayo \textit{The Methodology of Positive Economics} (1966) la prevalencia de la economía positiva, ya que considera que ésta será independiente de cualquier postura ética que se adopte. El único papel de la economía normativa para FRIEDMAN es supeditarla a los resultados de la economía positiva.\footnote{A modo de ejemplo, FRIEDMAN ciertamente lleva a cabo proposiciones normativas tales como su defensa de que las autoridades monetarias deben preocuparse por mantener controlado el ritmo de crecimiento de la oferta monetaria, pero esta aseveración se deriva de sus estudios de economía positiva que apuntan a la primacía de los incrementos de la oferta monetaria como factor determinante de la inflación.}; y, también,
        \item THALER, en su obra \textit{Towards a positive theory of consumer choice}, señala que la teoría de la elección tiene un carácter normativo.\footnote{THALER busca reorientar el enfoque del problema del consumidor desde un punto de vista de la economía positiva, es decir, respondiendo a la pregunta ¿cómo se comportan realmente los consumidores? Este autor argumenta que al trabajarse con seres racionales que maximizan su utilidad (preferencias egoístas que solo se preocupan por el beneficio propio), no se describe de manera fiel cómo es la realidad.}
    \end{lnum}
\end{la}

Más allá de este debate, es interesante recordar la nota discordante que puso JACOB VINER afirmando que \textit{economía es lo que hacen los economistas}, en una crítica velada a la delimitación de los contornos de la ciencia económica, y reflejo de que en la práctica, los economistas no prestan demasiada atención a lo que se supone debe ser el objeto de la economía y simplemente aplican sus métodos y su esfuerzo a las cuestiones que consideran merecen ser investigadas.\footnote{De todas formas, qué debe ser considerado o no economía sigue aún siendo objeto de debate. LIBERTAD GONZÁLEZ se defiende diciendo que  "\textit{(...) da igual como llamemos a mi trabajo, la pregunta es ¿este trabajo ayuda a entender algo?}". Por tanto, LIBERTAD ignora las cuestiones filosóficas y adopta una perspectiva más pragmática. Además, sus artículos empleaban econometría. Se han publicado en American Economic Review. \href{https://nadaesgratis.es/category/libertad-gonzalez}{https://nadaesgratis.es/category/libertad-gonzalez}.
}


\textcolor{red}{Además, aunque luego hablaremos de ello: El principio de organización hace referencia a la existencia de dos clases diferentes de teorías: las normativas y las descriptivas (o positivas). Las teorías normativas establecen la forma correcta de pensar sobre un problema, y por correcta no nos referimos en sentido moral, sino en sentido lógico y de coherencia. Un ejemplo paradigmático se encuentra en el corazón de nuestra ciencia: el modelo de optimización.}


\subsubsection{Ramificaciones: Microeconomía y/o Macroeconomía}
Podríamos realizar una primera clasificación en dos ramas: microeconomía y macroeconomía. Esta distinción comienza a hacerse más nítida tras la publicación de la Teoría General del Empleo, el Interés y el Dinero (1936) de JOHN MAYNARD KEYNES\footnote{\textit{Keynes’ General Theory, and his theory of unemployment, had an extraordinary impact. It changed the course of economic theory “by setting the scene for a new discipline, macroeconomics, that is simplified applied, and policy-oriented general equilibrium” (de Vroey, 2018)}.

MONACELLI, \textit{Macroeconomics and Economic Policy}: Lecture notes (2023)}. A raíz de esta obra se suele situar el origen de la macroeconomía como disciplina propia dentro de la ciencia económica.

\textbf{Disclaimer.-} La separación se ha complicado recientemente, pues la mayoría de la teoría macroeconómica desarrollada en la actualidad está microfundamentada. Por ello, hay controversias\footnote{\textit{Modern macroeconomics involves studying supply and demand in multiple markets at a time, rather than in a single one as microeconomists often do; but importantly where events in each market are allowed to depend on what is happening, and is expected to happen, in many or even all others. Thus it should be stressed that there are not different kinds of theories for microeconomics and macroeconomics. Any differences are fundamentally those of scope of the questions being asked (...).}

K. ATHREYA, \textit{Big Ideas in Macroeconomics} (2013)}.

Por ejemplo, el modelo de salarios de eficiencia de SHAPIRO y STIGLITZ (1984) aparece en manuales de macroeconomía ya que supone una microfundamentación de rigideces en el mercado de trabajo (algo fundamental en la macroeconomía de la NEK) y una explicación a una variable agregada (desempleo). Sin embargo, en puridad, el análisis del mercado de trabajo que realizan es en equilibrio parcial y, por tanto, se puede argumentar que es un análisis microeconómico.

\paragraph{Microeconomía}
La microeconomía es la rama de la ciencia económica que se encarga del estudio individualizado del funcionamiento de los mercados. Serán los economistas neoclásicos quienes se centren en el desarrollo de la teoría microeconómica, sentando los cimientos del paradigma científico actual.\footnote{La idea clave que hace que prolifere este análisis es el empleo de la teoría de la elección aplicada sobre los agentes económicos, que hace uso del cálculo diferencial. Así, a finales del siglo XIX y principios del siglo XX, estos autores sientan los cimientos de la teoría de la demanda, las teorías de la producción y de costes, y el análisis de las estructuras de mercado [\authorfont{Ver Temas 3.A.8-3.A.10, 3.A.11-3.A.13 y 3.A.15-3.A.21}, respectivamente]. En este sentido, cabe destacar el papel de LÉON WALRAS y ALFRED MARSHALL}

A partir de la teoría neoclásica se desarrollan una serie de campos que parten de la teoría de la elección:
\begin{lnum}
    \item Economía del bienestar: Vertiente normativa del estudio de los mercados. Destacan las contribuciones de PARETO, PIGOU y COASE [ver temas 3.A.22, 3.A.23 y 3.A.24].
    \item \textit{Economía pública.} Donde destacan las aportaciones de la teoría de la imposición óptima de RAMSEY, MIRRLEES y ATKINSON.
    \item Economía de la información: Desarrollada a partir de la década de 1970 con las aportaciones de AKERLOF [ver tema 3.A.13]. 
    \item \textit{Economía del mercado de trabajo.} Estudiado por distintos autores, desde MARSHALL y PIGOU o HICKS hasta BECKER [ver temas 3.A.25-3.A.27].
    \item \textit{Comercio internacional.} Busca estudiar las causas del comercio entre países, el patrón de comercio y los costes y beneficios del comercio internacional. Destacan las aportaciones de OHLIN, SAMUELSON, KRUGMAN y MELITZ [ver temas 3.B.5-3.B.9].
    \item \textit{Economía financiera.} Destaca el análisis seminal de MARKOWITZ (modelo de media-varianza) que es una aplicación de la teoría de la elección en situaciones de riesgo [ver temas 3.B.23-3.B.25].
    \item \textit{Economía del comportamiento.} Incorpora ideas y resultados de la psicología y la sociología en la modelización teórica de diversas cuestiones económicas. Es un claro ejemplo de interdisciplinariedad, la otra cara de la moneda del imperialismo económico. En este campo destacan las aportaciones de KAHNEMAN y THALER.
    \item Otras disciplinas, en parte estimuladas por BECKER, como Economía de la salud o Economía de la educación.
\end{lnum}

\paragraph{Macroeconomía}
Por su parte, la macroeconomía es la disciplina que estudia el funcionamiento de la economía en su conjunto, valiéndose de la agregación de la conducta de los agentes. A partir de la crítica de LUCAS, todas las ramas de la macroeconomía se basan en modelos dinámicos microfundamentados:
\begin{lnum}
    \item Teoría del crecimiento económico [\authorfont{Ver Temas 3.A.43-3.A.45}];
    \item Teoría de ciclos económicos [\authorfont{Ver Tema 3.A.42}].
    \item Economía internacional (modelos de la NOEM) [\authorfont{Ver Temas 3.B.15-3.B.16}].
\end{lnum}

Desde este punto de vista, se puede argumentar que, a nivel metodológico, la macroeconomía moderna le debe más a IRVING FISHER que a KEYNES [\authorfont{Ver Tema 3.A.29}]. Sin embargo, la influencia de KEYNES ha sido decisiva en la macroeconomía actual, al menos en implicaciones de políticas económicas y concepción de rigideces (por ejemplo, NEK).\footnote{KEYNES da origen a la macroeconomía como rama independiente, pues su obra es un claro ejemplo de un fenómeno recurrente de la teoría económica: la situación social influye sobre el objeto estudiado. Sin embargo, aseverar que KEYNES aporta los cimientos de la teoría macroeconómica actual es más aventurado dado el desarrollo de la macroeconomía hacia la microfundamentación con modelos eminentemente dinámicos que ocurre a partir de las contribuciones de ROBERT LUCAS emergencia de la Nueva Macroeconomía Clásica.}




\subsection{Enfoque: Parcial y/o General} 
En economía:
\begin{la}
    \item Un modelo de equilibrio general es aquel cuyas predicciones resultan de tener en cuenta todos los parámetros, interacciones, decisiones, reglas, contexto institucional, etc. que se consideran relevantes para representar el fenómeno en cuestión;\footnote{ Los modelos de equilibrio general representados formalmente en lenguaje matemático tienen su origen en WALRAS (1873).} y, por otro lado,
    \item Un modelo de equilibrio parcial es aquel que reconoce la existencia de una serie de factores como relevantes para dar lugar a un fenómeno determinado, pero en el que conscientemente se decide obviar uno o varios de esos factores para obtener una formulación más tratable y que logre transmitir de manera más efectiva una determinada explicación de un fenómeno económico. Esta decisión de obviar determinados factores se materializa generalmente en la aplicación del supuesto de \textit{ceteris paribus}.\footnote{La obra de ALFRED MARSHALL es habitualmente considerada como un ejemplo paradigmático del enfoque de equilibrio parcial, dado el uso generalizado del supuesto de \textit{ceteris paribus}.
    
    Aunque en ocasiones se contrapone el enfoque de ambos autores como si se hubiese tratado de un sujeto de controversia entre ambos, lo cierto es que ambos enfoques son complementarios, y ni siquiera es probable que ninguno de ambos conociese la obra del otro lo suficiente como para elaborar una teoría deliberadamente antagónica.}
\end{la}  

En la práctica, las fronteras entre equilibrio general y parcial son a menudo difusas. Dados los inevitables límites computacionales e informacionales que existen en toda representación de un fenómeno económico, resulta imposible incluir todos los supuestos y factores relevantes a este. Ello resulta en que los modelos definidos como de equilibrio general incluyen simplificaciones de fenómenos que son evidentemente relevantes, pero que se obvian para mantener la tratabilidad del modelo.\footnote{El principal exponente actual de los modelos de equilibrio general, los denominados modelos DSGE\footnote{Dynamc Stochastic General Equilibrium} (en todas sus variantes), obvian fenómenos e interacciones que se reconocen como relevantes pero que se consideran intratables o no lo suficientemente relevantes para ser incluidos en la formulación del modelo.}


\subsection{Controversias: El dinero}
Si bien es cierto que muchos problemas de elección que enfrenta el agente no suscitan controversia, hay un elemento particular, en el corazón del funcionamiento de la economía, sobre el cual los académicos tienen por costumbre debatir con frecuencia: el dinero. ¿Qué papel juega? ¿Forma parte de las alternativas de elección?

Por otro lado, determinar la neutralidad o no del dinero constituye un elemento crucial a la hora de delimitar el objeto de la ciencia económica. Para entender el por qué solo es necesario hacerse una pregunta, ¿es necesario incluir dentro del objeto de estudio un elemento sin influencia ni interacción en el outcome estudiado? Para entender a lo que nos referimos, pensemos si es o no necesario incluir las dinámicas solares, las tormentas, emisiones, etc. dentro de un modelo económica. Aunque constituyen una parte evidente de la realidad en la que habitamos, ¿debemos estudiarlas dentro del campo de nuestra ciencia? Evidentemente no. 

No obstante, en el desarrollo de la teoría económica las reflexiones sobre el rol del dinero siempre han generado un espacio de especial controversia. Este tema ha sido objeto de intensos debates en los últimos cinco siglos: \textit{La importancia del dinero en el corto y en el largo plazo es uno de los ejes que dividen el pensamiento económico}. 

\textcolor{red}{\textbf{OJO!} A partir de aquí todo puede no tener sentido, la última modificación fue solo de reestructuración}


Hay quienes piensan que el dinero es solo un velo, totalmente neutral al quehacer económico, y hay quienes aseguran que el dinero cumple un rol central, al menos en el corto plazo. La pregunta será entonces,

\textit{¿Afecta el dinero la actividad económica real? Y si la afecta, ¿cuál es el mecanismo de transmisión de cambios en el acervo de dinero al sector real de la economía?} 

La pregunta central que se hace la teoría es sobre la conexión que existe entre los fenómenos monetarios y los fenómenos reales de una economía, es decir si el dinero cumple un rol más allá de la economía de trueque. 

\subsubsection{No-Neutralidad del dinero: KEYNES}
En sus Ensayos monetarios, KEYNES refutó la TQM argumentando que ante un cambio en la cantidad de dinero, no todos los precios varían en la misma proporción. Para KEYNES el dinero cumple un rol central en la economía y planificación, ya que existen cambios en los precios relativos y por tanto distintos efectos distributivos con sus consiguientes impactos en la producción, pudiendo el dinero impulsar la inversión y ésta transformarse en motor del crecimiento.\footnote{Esta idea rompe la causalidad directa de la oferta monetaria sobre los precios y da al dinero un rol central. 

No obstante, KEYNES reconoce que en el largo plazo la economía tiende a funcionar como señalan los clásicos, aunque el largo plazo es un período de tiempo irrelevante dado que en el largo plazo \textit{todos estaremos muertos}.} Además, para éste el desempleo es persistente y solo con políticas de demanda agregada que impulsen la demanda efectiva se puede alcanzar el equilibrio sin desempleo involuntario.\footnote{La existencia de este espacio para que la política monetaria y fiscal pueda mejorar el nivel de ingreso, es para KEYNES la prueba irrefutable de que la neutralidad del dinero es una falacia.}

En \textit{Teoría General del empleo, el interés y el dinero}, KEYNES se plantea formalmente la vinculación entre las variables reales y las variables monetarias del sistema a través del efecto de los precios. El mecanismo de transmisión y de fluctuación en el valor de los activos financieros es por la vía de la tasa de interés. De este modo, el dinero si cumple un rol central en la economía, dado que junto a su desempeño como medio de cambio (economía de trueque) y atesoramiento, añade su rol como vehículo para la especulación. Con esto, la antigua determinación monetaria de los precios es suplantada por el argumento keynesiano de la demanda efectiva y la tasa de interés. Por tanto, no es un velo neutral sino que es el elemento que une el presente con el futuro, que aumenta la brecha entre la esfera real y la esfera monetaria debido al pago de intereses que, en puridad, se trata de transferencias de dinero hacia los dueños del capital, y acaba potenciando la desigualdad.\footnote{Si bien los intereses son la expresión de que en el capitalismo las decisiones económicas se toman en condiciones de incertidumbre, ante la incertidumbre del futuro muchas personas atesoran dinero, y este atesoramiento de dinero por parte de las empresas y las familias puede afectar el ciclo económico.} 

El consenso internacional alrededor de la obra de KEYNES fue abrumador y llego a ser el paradigma científica dominante entre los años 30's y los años 70's.\footnote{Este ciclo progresista, que se proyecta entre los años 30 hasta mediados de los años 70 (cuarenta años), marcado por el New Deal y la influencia de la obra de éste autor, es probablemente el período de mayor bonanza económica en la historia humana. PAUL KRUGMAN ha demostrado que en este período mejoró la distribución del ingreso no sólo en los Estados Unidos sino en todas las economías emergentes.¿La razón?: el desarrollo de políticas encauzadas al pleno empleo.
La políticas que velaron por el pleno empleo a la larga resultaron más eficaces que aquellas que sólo miraban el control de precios. Este único parámetro que imponían a tabla-rasa los bancos centrales, el FMI y el ya alcaído consenso de Washington quedará relegado a la nada. Fuente: \href{https://www.elblogsalmon.com/historia-de-la-economia/el-adios-a-la-contrarrevolucion-monetarista}{El adiós a la contrarrevolución monetarista de Milton Friedman}}


\subsubsection{Neutralidad del dinero: Monetarismo}
La respuesta de toda la economía clásica, neoclásica y monetarista es siempre la misma: \textit{el dinero es neutral, es solo un velo del sistema}, y no influye en la economía ni en el corto ni en el largo plazo.\footnote{La idea sobre la neutralidad del dinero se atribuye con frecuencia a las obras de DAVID HUME (\textit{Of Money}, 1750), aunque fue FRIEDRICH VON HAYEK quien la conceptualizó en su obra \textit{Price and Production} (1931). Sin embargo, cuando HUME teorizó acerca de la neutralidad del dinero y, en particular, sobre la Teoría Cuantitativa del Dinero, esta idea tenía más de 200 años de antiguedad cuando el filósofo de \textit{La naturaleza Humana} y maestro de ADAM SMITH la desarrolla en \textit{Of Money}. 

El filósofo y político francés JEAN BODIN la formuló en 1568 (182 años antes que HUME) en su primer texto de teoría monetaria. En esta obra, BODIN señala que la principal causa del aumento de los precios es el aumento de las cantidades existentes de oro y plata.

Esta idea tuvo gran impacto en Europa y fue considerada por largo tiempo como la primera exposición de una teoría cuantitativa del dinero. Sin embargo más tarde se descubrió que la primera construcción científica sobre el tema fue elaborada por los pensadores de la Escuela de Salamanca, en concreto por el sacerdote jesuita MARTIN DE AZPILCUETA, en 1525 (225 años antes de HUME). AZPILCUETA describió los efectos inflacionistas de la masiva llegada de oro y plata producto del saqueo realizado por los conquistadores y navegantes a los confines del mundo desde el siglo XIV.} Es decir, las variaciones exógenas en la oferta monetaria solo afectan en igual proporción a las variables nominales, dejando inalteradas a las variables reales como la producción, el empleo y los precios relativos. De acuerdo con estas escuelas, los fenómenos monetarios no afectan los fenómenos reales de la economía y solo tienen impacto sobre las nominales, en particular, el nivel de precios. 

La Teoría Cuantitativa del Dinero (TQM) ha constituido el núcleo central y el paradigma macroeconómico dominante utilizado por numerosos economistas que contribuyeron a su desarrollo, como WICKSELL, FISHER, MARSHALL, PIGOU, y HAYEK, entre otros.\footnote{Asociada estrechamente a los economistas clásicos y a la Escuela de Chicago} Estos economistas consideran que el objetivo principal de la política económica es conservar el valor del dinero.

La formulación de la Teoría Cuantitativa es atribuible a IRVING FISHER (en su versión más moderna) en su obra \textit{The Purchasing power of Money}, 1911, en la que expone lo que se ha dado en llamar el enfoque de transacciones de la teoria cuantitativa: 

$$MV=PT$$

Esta identidad expresa que la cantidad de dinero existente en la economía ($M$), por el número de veces que se usa cada unidad monetaria ($V$) es equivalente al valor total de las transacciones efectuadas ($P\;\cdot\;T$). De acuerdo a esta relación, la causalidad va de izquierda a derecha y es la cantidad de dinero la que determina el nivel de precios.\footnote{De acuerdo a esta identidad de FISHER, y siempre bajo el supuesto de una economía en pleno empleo (regla general de los economistas clásicos), la velocidad de circulación del dinero ($V$) es constante por razones tecno-institucionales determinadas por los hábitos de pago y de gastos de la comunidad. Del mismo modo, el nivel de actividad o número de transacciones ($T$) se asume constante, lo que deja el nivel de precios ($P$) en función de la cantidad de dinero ($M$).}

\textsc{Nuevos desarrollos}

Dado que el enfoque de FISHER incluye el total de transacciones (finales e intermedias) y esto plantea el problema contable de duplicidad de operaciones, surgió el enfoque de ingresos de la teoría cuantitativa, llamada también ecuación de cambio: 

$$MV=PY$$

En esta, se considera que el valor del ingreso real de la economíoa ($P\;\cdot\;Y$) es equivalente a la demanda monetaria ($M$), por el número de veces que, en promedio, se usa cada unidad monetaria en operaciones económicas ($V$). Si bien esta versión es más operativa, sigue siendo una identidad similar a las anteriores que expresa una causalidad entre cantidad de dinero y nivel de precios, con causalidad de izquierda a derecha. Además, para FISHER la Velocidad del Dinero es constante pero, comoo vemos, esto contrasta con la evidencia empírica:

\imagenfit{VelM2StockUS.png}{Velocity of M2 Money Stock (Recesiones: areas grises)}{Federal Reserve Bank of St. Louis}

Está claro que si el dinero es neutral, como señalan los clásicos, neoclásicos y monetaristas, no tiene sentido que el gobierno intervenga a través de la política monetaria, dado que el banco central no puede corregir el sistema aumentando o disminuyendo la cantidad de dinero para suavizar las fluctuaciones del producto y el empleo en torno a sus niveles de equilibrio.


La contrarrevolución monetarista inició su desarrollo durante los años 50's contra las ideas de KEYNES, intentado demostrar la neutralidad del dinero a través de la tesis del equilibrio general walrasiano y uno de sus grandes exponentes fue MILTON FRIEDMAN.\footnote{Desde que LEON WALRAS, el fundador de esta teoría (1874), modeló el concepto de equilibrio para una economía con dinero, numerosos autores han perfeccionado el modelo matemáticamente al crear la heurística positiva que lo retroalimenta. WALRAS, en su sistema de ecuaciones de excesos de demanda y excesos de oferta de mercancias que se ajustan automáticamente hasta vaciar el mercado por la vía de los precios, introdujo una ecuación que representa los excesos de demanda y oferta de dinero. En su sistema de infinitas ecuaciones, WALRAS no se percató que el equilibrio del sistema se establece con muchos precios negativos, cosa que es imposible en la vida real. Pese a que los modernos sistemas computacionales evitan este problema, lo importante en la idea de WALRAS es que advierte que el dinero es algo que carece de utilidad para los que lo poseen, y que su servicio sólo satisface la necesidad de transacciones (trueque).

La concepción del dinero en el equilibrio general walrasiano fue conceptualizada en los años 50 en el Modelo ARROW-DEBREU, gracias a los aportes matemáticos realizados en los años 30's por WALD y VON NEUMANN. El documento de KENNETH ARROW y GERARD DEBREU \textit{Existence of an Equilibrium for a Competitive Economy}, 1954, es un modelo matemático de alta complejidad que introduce la noción de incertidumbre bajo la idea de localización: una manzana no vale lo mismo en Europa o en África; ni tampoco vale o mismo hoy o en 30 días. La incertidumbre de los precios diferentes en una u otra parte del mundo, y en el presente o futuro, no es más que un agregado de ecuaciones en las que lo medular del rol del dinero no salta a la vista.

En el modelo del equilibrio general de ARROW-DEBREU el dinero sigue siendo neutral y no desempeña un rol determinante en la economía.} 

El hecho que permitió la arremetida de FRIEDMAN fue la crisis petrolera de los años 70's, que elevó los precios y llevó la inflación a niveles insostenibles en algunos países. Para FRIEDMAN, la inflación era consecuencia de las políticas de pleno empleo y de los gastos desmedidos de los gobiernos que obligaban a imprimir más dinero: \textit{la inflación es siempre y en todo momento un fenómeno monetario}, como enfatizó MILTON FRIEDMAN..

Por eso que los gobiernos tenían que limitarse a mantener sólido el dinero y dejar que la economía se cuidara a sí misma.\footnote{Por eso que la única intervención defendida por éstos es el control de la inflación, en la sospecha de que éste es un fenómeno estrictamente monetario (FRIEDMAN). Los clásicos y neoclásicos coinciden en que la política monetaria debe ajustar la cantidad de dinero de tal manera que conduzca a la estabilidad de precios; pero hay precios que escapan a la estabilidad y no hace más que confirmar la dicotomía clásica en torno al rol del dinero y por qué los fenómenos reales no logran dar cuenta de los fenómenos monetarios de la economía. La tasa de interés real fue el mecanismo de unión entre la esfera monetaria y la esfera real: (i) la manipulación de los tipos de interés; y, (ii) la economía fiduciaria generada por los bancos privados; hace aún más complejo los mecanismos de transmisión del dinero de la esfera monetaria a la esfera real. Por ello, cabe incluso destacar, que la teoría económica ha sido incapaz de introducir a los mercados financieros en sus modelos de equilibrio general (nada negligible pues sólo los mercados de derivados financieros equivalen a más de diez veces la producción económica global).} La contrarrevolución monetarista, como se dio a conocer, postuló que sacar al gobierno de la actividad económica permitiría que las fuerzas del mercado (al igual que las fuerzas gravitatorias) convergieran naturalmente hacia el pleno empleo con tasas de crecimiento cada vez más elevadas y con una distribución (vía teoría del chorreo) cada vez más eficiente.\footnote{La promesa de esta contrarrevolución monetarista hizo agua por todos lados: no sólo nunca logró el pleno empleo sino que generó uno de los períodos de mayor desigualdad económica. El mito de la autorregulación de los mercados quedó develado y se demostró que las economías de mercado no son naturalmente estables sino que deben ser estabilizadas mediante políticas.}

KEYNES hizo hincapié en la fragilidad de las expectativas sobre las que se basa la actividad económica en los mercados descentralizados: el futuro es incierto y la psicología del mercado es caprichosa.


\imagenfit{MonestaryBaseFED_USA.png}{St. Louis Adjusted Moneatry Base (Recesiones, areas grises)}{Federal Reserve Bank of St. Louis}


Como hemos visto, los supuestos de la teoría neoclásica para que el dinero sea neutral son altamente restrictivos, lo que convierte este concepto en un interesante ejercicio teóríco (en la línea de la Hipótesis de los mercados eficientes, de Eugene Fama), que no pasa de ser una elucubración matemática que oculta el verdadero significado del dinero y, más aún, en la creación de dinero desde la nada.



El debate sobre la neutralidad del dinero gira en torno a si los cambios en la oferta de dinero afectan solo a los precios (neutralidad) o también a variables reales como el empleo y la producción (no neutralidad). Las escuelas clásicas y monetaristas argumentan que el dinero es neutral a largo plazo, aunque puede tener efectos en el corto plazo. Por otro lado, las escuelas keynesianas y austriacas sostienen que el dinero no es neutral, ni siquiera a largo plazo, debido a la rigidez de precios y expectativas. 

Por otro lado, La neutralidad del dinero es la idea de que un cambio en la cantidad de dinero afecta sólo a las variables nominales de la economía, tales como precios, salarios y tipo de cambio, sin afectar a las variables reales, como el empleo, el producto interno bruto real o el consumo.\footnote{Patinkin, Don (1987), \textit{Neutrality of Money}}

Acotaremos el debate a sólo dos de estas escuelas,\footnote{RAVIER (2010), \textit{La no-neutralidad del dinero a largo plazo. Un debate entre Chicago y Viena}}: la Escuela de Chicago o monetarista y a la Escuela Austríaca. Sobre éstas podría decirse que sintetizan la existencia de dos tradiciones claramente diferenciadas con respecto al estudio de los efectos que la política monetaria ejerce sobre la actividad económica y el empleo:
\begin{la}
    \item Los monetaristas ortodoxos (Escuela de Chicago) se concentran sobre los cambios en el nivel general de precios que la expansión y la contracción monetaria traen aparejadas; todos los precios, se asume, se mueven hacia arriba o hacia abajo uniformemente. Esto se mantiene en parte porque las magnitudes como el nivel general de precios son fácilmente observables, y en parte porque se asume siempre que el dinero es neutral (Barry, 1981, 23).\footnote{La Escuela de Chicago se presenta como continuadora del tratamiento agregado representado en John Locke, David Hume y la mayoría de los pensadores clásicos.}
    \item En contraste, la tradición austríaca no niega el efecto no neutral de la política monetaria o los efectos del shock monetario sobre los precios rela- tivos. Los austríacos enfatizan estos efectos y la resultante distorsión de la actividad real.\footnote{La Escuela Austríaca es la heredera del tratamiento desagregado iniciado en Richard Cantillon (1755) y seguido por, quizás el último de los clásicos, John Cairnes (1873) (Ravier, 2008b).}
\end{la}  


La neutralidad del dinero es una idea importante en la economía clásica y se relaciona con la dicotomía clásica. Esto implica que el banco central no afecta a la economía real (por ejemplo, el número de puestos de trabajo, el tamaño del PIB real, el monto de la inversión real) mediante la emisión de dinero. En cambio, cualquier aumento en la oferta de dinero se vería compensado por un aumento proporcional de los precios y los salarios. Este supuesto refuerza algunos modelos macroeconómicos convencionales (por ejemplo, la teoría de los ciclos económico reales) mientras que otros como el monetarismo ven al dinero como neutral solamente en el largo plazo.

\paragraph{Superneutralidad del dinero: NMC}
La superneutralidad de dinero es una propiedad más fuerte que la neutralidad del dinero. Sostiene que no sólo la economía real no es afectada por el nivel de la oferta de dinero, sino también que la tasa de crecimiento de la oferta monetaria no tiene ningún efecto sobre las variables reales. En este caso, los salarios nominales y los precios siguen siendo proporcionales a la oferta nominal de dinero, no sólo en respuesta a los cambios permanentes únicos en la oferta nominal de dinero, pero también en respuesta a los cambios permanentes en la tasa de crecimiento de la oferta monetaria nominal. Normalmente la superneutralidad se aborda en el contexto de los modelos de largo plazo.



\clearpage
\section{Ciencia Económica: Método}
\puntosclave{
    La epistemología es el estudio de las diferentes formas de alcanzar nuevo conocimiento, incluyendo el conocimiento económico. 
    
    La deducción y la inducción caracterizan las dos vías esenciales para obtener nuevo conocimiento, y ambas son relevantes en economía.
    
    El falsacionismo es una corriente epistemológica que aboga por exigir que las teorías científicas puedan ser probadas falsos para considerarse útiles a la obtención de nuevo conocimiento. Algunas teorías relevantes en economía han sido criticadas por su difícil falsabilidad.
    
    Los paradigmas científicos y los programas de investigación son herramientas epistemológicas e historiográficas que tienen por objetivo explicar la aparición de nuevas teorías y métodos que expliquen fenómenos no explicados de manera satisfactoria por las teorías existentes. 
    
    El estudio de los fenómenos económicos ha dado lugar a una serie de controversias persistentes: ¿es importante que los supuestos de un modelo sean realistas?, ¿la economía debe estudiar el qué es o el qué debe ser?
    
    Muchos de los fenómenos que estudian los economistas no tienen una explicación satisfactoria, y aparecen nuevas técnicas metodológicas con el objetivo de superar tales explicaciones insatisfactorias.
}

La economía. en tanto que conjunto de teorías y métodos para entender y predecir una serie de a fenómenos, es susceptible de valorarse en los mismos términos que otras teorías científicas. Así, cabe recorrer brevemente la influencia de los cambios en el método científico y la epistemología de la ciencia sobre la economía (o la ciencia económica).

\subsection{Epistemología: Breve introducción}
Se denomina epistemología a la rama de la filosofía que trata de caracterizar los diferentes métodos de obtención de nuevo conocimiento, atendiendo, principalmente a los determinantes lógicos, psicológicos, políticos, computacionales e históricos, entre otros, que influyen en ese proceso. El método científico es un conjunto de técnicas utilizadas para investigar las causas de fenómenos observados, adquirir nuevos conocimientos e integrar el conocimiento obtenido con otras explicaciones de esos fenómenos.\footnote{DESCARTES señaló que el método es necesario para la investigación de la realidad. Este camino consta a su vez de dos fases:
\begin{lnum}
  \item Heurística. Elaboración de una teoría que dé cuenta de los hechos o fenómenos observados.
  \item Dialéctica. Contrastación de las proposiciones teóricas con las correspondientes fuentes informativas originales.
\end{lnum}
} 

Los medios con los que el creciente mundo científico trató de descubrir la verdad implicaban la observación empírica cuyo principal ejemplo es el método científico, que entrañaba la integración de la razón y la observación empírica. Aunque este tema es demasiado complicado para explicarlo más detalladamente,\footnote{¿Cómo hacemos para responder a las preguntas \texit{qué sabemos} y \textit{cómo sabemos que lo que sabemos es correcto}? Depende de la respuesta a la pregunta \textit{¿existe una verdad última que los científicos estén tratando de revelar (postura absolutista) o no existe ninguna verdad subyacente (postura relativista)?}
\begin{la}
  \item Si existe una verdad última, ¿cómo la encontraremos?
  \item Si no existe ninguna, ¿hay algunas proposiciones más verdaderas que otras?
\end{la}
Los spistemólogos pasados y presentes no han sido capaces de ponerse de acuerdo sobre estos problemas, pero han generado abundante literatura sobre el tema. Si creemos que existe una verdad última, el problema es cómo saber que ya la hemos descubierto.} la verificación se analiza minuciosamente en los escritos de KANT, HUME, DESCARTES y otros filósofos de los siglos XVII y XVIII.

\subsubsection{Método deductivo: DESCARTES}
La deducción consiste en la derivación lógica de una o varias  conclusiones a partir de unas premisas que se asumen ciertas. Es  habitual resumir este método como un recorrido \textit{de lo general  a lo particular},\footnote{El método hipotético-deductivo tiene su origen en el Discurso del Método (1637) de DESCARTES:
\begin{lnum}
  \item Se plantea un conjunto axiomático de partida. Debe tenerse en cuenta que el concepto de axioma hace referencia a una verdad evidente por sí misma. Por tanto, un axioma no es equivalente a un supuesto;
  \item Partiendo de los axiomas observados, se lleva a cabo un proceso de deducción lógica con el fin de obtener una serie de conclusiones; y, por último,
  \item Finalmente, a partir de las conclusiones, se intentarán deducir leyes de carácter general que se enunciarán en forma de teoremas que al venir de axiomas serán verdades evidentes.
\end{lnum}
Los defensores de este método se preocupan por la relevancia de sus conclusiones y no por la validez que viene dada por definición.

Sin embargo, cabe destacar la dificultad de partir de axiomas que no generen disputa por su condición de irrefutables.} en el sentido de que mediante la aplicación conjunta  de una serie de axiomas es posible establecer un resultado, una predicción o una conclusión para la situación particular que habían definido esos axiomas generales. 

En economía, el método deductivo ha sido utilizado profusamente en economía:
\begin{lnum}
  \item RICARDO fue pionero en la construcción de modelos económicos, herramienta básica en la metodología económica actual;\footnote{Un modelo es una construcción ficticia que busca ser un molde para explicar y dar coherencia a una teoría. Se busca priorizar un determinado aspecto de la teoría (i.e. se realiza un ejercicio de abstracción, ya que los fenómenos económicos son complejos, por lo que se quiere evitar detalles y complicaciones accesorias e innecesarias).  Un modelo está formado por un conjunto de supuestos, a partir de las cuales se derivan, mediante las correspondientes pruebas matemáticas (incorpora cierto grado de formalismo), algunas conclusiones. Por eso, en este sentido, un modelo sería una construcción análoga a un sistema deductivo.} y,
  \item J.S. MILL se decanta principalmente por el método deductivo.\footnote{Una aportación metodológica fundamental, fruto de la construcción axiomática propia del método hipotético-deductivo, es el concepto de \textit{homo economicus}, que conduce a suponer que los individuos son seres racionales que se guían por incentivos económicos que buscan lograr su máximo bienestar.}
\end{lnum}

Sin embargo, eesta metodología fue criticada por lo que SCHUMPETER denominó vicio ricardiano (\textit{ricardian vice}), consistente en la creación de modelos abstractos con supuestos poco realistas y sin verificación empírica.\footnote{A pesar de esto, desde el desarrollo de los modelos micro y macroeconómicos basados en representaciones formales, la deducción es una herramienta esencial para la hipotetización de nuevas teorías, cuya capacidad predictiva y consistencia empírica trata de contrastarse posteriormente mediante la comparación con observaciones reales cuantificadas y analizadas mediante técnicas estadísticas y herramientas informáticas.
}

\subsubsection{Método inductivo: HUME}
La inducción o método inductivo toma el camino opuesto y, de una serie de observaciones sobre fenómenos particulares, trata de derivar un resultado general.\footnote{Esta metodología se asocia originalmente a los trabajos de BACON a comienzos del siglo XVII. En términos generales, podría definirse como un método que consiste en establecer enunciados universales ciertos a partir de la experiencia. Así, las investigaciones científicas basadas en el método inductivo seguirían el siguiente proceso:
\begin{lnum}
  \item Observación de los hechos de forma libre y carente de prejuicios, con el fin de descubrir regularidades;
  \item Con posterioridad, y mediante inferencia, se formulan leyes universales en base a los hechos observados; y,
  Finalmente, por inducción, se obtienen afirmaciones más generales que recibirán el nombre de teorías.
\end{lnum}

Según los empiristas clásicos, sólo pueden considerarse teorías científicas aquellas que estén formadas por un conjunto de enunciados probados empíricamente. Esto implica que: (i) una teoría no es aceptada hasta que no haya sido probada; (ii) las teorías pueden ser verificadas a través de la contrastación empírica; y, (iii) rechazo frontal hacia toda especulación teórica para la cual no pueda realizarse una contrastación empírica. Por tanto, se propugna el establecimiento de enunciados universales obtenidos a partir de la experiencia conocidos como teorías, diferentes de los teoremas al provenir de la experiencia. Por tanto, sólo podrán considerarse teorías científicas aquellas que vienen de la experiencia. 

Si, por ejemplo, tomamos una población de empresas durante tres años, y vemos que aquellas que tuvieron un beneficio operativo positivo en ese periodo seguían activas en el periodo siguiente, podríamos inducir que beneficio operativo positivo es una condición suficiente para la supervivencia de una empresa. Sin embargo, nada en ese proceso inductivo implica que en el cuarto año no tendrá lugar un endurecimiento de las condiciones de financiación que reduzca el beneficio neto de las empresas que estaban endeudadas, y con ello se produzcan ceses de actividad en empresas que efectivamente seguían teniendo un beneficio operativo positivo. Es precisamente ésta la principal crítica al método inductivo: ¿es acaso posible garantizar que no aparecerá una observación incompatible con la relación inducida por las observaciones previamente conocidas?\\
Este problema se conoce como \textit{problema de la inducción}, y radica en si un resultado obtenido mediante inducción está justificado epistemológicamente, es decir, si la inducción produce conocimiento. Fue introducido por DAVID HUME [\authorfont{Ver Anexo \ref{anx:espitemologia}}]}

En economía, la inducción ha sido utilizada de manera habitual y ya desde el periodo previo a los clásicos. Algunos ejemplos de método inductivo en economía son:
\begin{lnum}
  \item Esta concepción empírica de la ciencia está perfectamente resumida en la obra de J.S. MILL, \textit{System of Logic, Ratiocinative and Inductive} (1843), aunque buena parte de sus desarrollos económicos no fueron coherentes con sus postulados metodológicos;
  \item Preferencia revelada de SAMUELSON: Se parte de la cesta consumida por el individuo en el mercado para buscar establecer una serie de conclusiones y generalidades en base a dicho estudio. En cambio, la teoría ordinal de la demanda del consumidor sigue un enfoque deductivo: se parte de la formalización de unas preferencias (axiomas), se resuelve el problema del consumidor y se obtiene la demanda [\authorfont{Ver Tema 3.A.8}];
  \item Escuela de Harvard de Organización Industrial: Parte de la observación empírica (en Estados Unidos se están imponiendo grandes empresas sobre pequeñas) [\authorfont{Ver Tema 3.A.15}]; y,
  \item En la actualidad, el desarrollo de técnicas estadísticas y las posibilidades crecientes de recogida de datos y procesamiento de información han hecho de la  econometría la principal herramienta de aplicación de la inducción para expandir el conocimiento económico,\footnote{Ejemplo de esto último es la reciente proliferación del uso de experimentos en economía, tanto naturales como controlados. Recordar, Premio Nobel de Economía 2021 a CARD, ANGRIST y IMBENS, "\textit{Por sus contribuciones metodológicas al análisis de las relaciones causales}"} por ejemplo, Randomized Control Trials (RCTs) [\authorfont{Ver Tema 3.B.28}].
\end{lnum}

También se enfrenta a algunos problemas particulares, entre otros:
\begin{lnum}
  \item No hay garantía de que la teoría verificada en $t$ se verifique en $t+1$ (problema de la inducción de HUME);
  \item Imposibilidad de recopilar todos los hechos relacionados con los fenómenos de estudio; y,
  \item El método inductivo no puede garantizar que las generalizaciones que defienda a partir de la observación sean válidas.
\end{lnum}

\subsubsection{Síntesis: KANT}
En su obra \textit{Crítica de la razón pura} (1781), KANT aboga por la complementariedad de los métodos. Según este filósofo, tanto la razón como la experiencia juegan un papel importante en el conocimiento. Concretamente, la razón humana (asociada al deductivismo) tendrá un carácter a priori en la construcción del conocimiento, mientras que la experiencia (más relacionada con el inductivismo) tendrá un carácter a posteriori. Dicho de otro modo: \textit{es imposible demostrar nada que no haya sido previamente enunciado}. Como la mayoría de los filósofos cree que el conocimiento se deriva de una combinación de los dos. 

De heho, el debate girará en torno a la naturaleza de la combinación óptima. Por eso, cabe tener en cuenta, que ni inducción ni deducción se aplican de manera perfectamente estanca en la práctica de la economía. Los modelos econométricos aplican resultados matemáticos bajo los cuáles subyacen inevitablemente herramientas matemáticas en las que la deducción juega un papel central. De manera similar, modelos matemáticos tales como el Modelo del Ciclo Real, el modelo neoclasico de la decisión del consumidor o los incontables otros modelos formulados en el último siglo en lenguaje matemático se formularon motivados por la observación previa de una serie de fenómenos particulares respecto de los cuales cabía inducir una regla general.

\subsection{Controversias sobre el método: Methodenstreit}
La segunda mitad del siglo XIX vendrá marcada en economía por la coexistencia de dos escuelas que difieren sustancialmente en su metodología científica:\footnote{La batalla del método continúa en el siglo XX, enfrentando a ROBBINS con HUTCHINSON. ROBBINS (1932) defiende el método deductivo. El cuerpo central de las proposiciones de la teoría económica ortodoxa se basaba en proposiciones que eran tautologías, lo que conducía con frecuencia a una posición circular, en cuanto se suponía como axiomático lo que era necesario probar. Por añadidura, HUTCHISON alerta que el interés en los problemas de política económica ha declinado significativamente, forma de un rigor matemático vacío de contenido empírico.}
\begin{la}
  \item \textit{Historicismo alemán.} Los autores más representativos son SCHMOLLER, HILDEBRAND y LIST, quienes rechazan que en economía se puedan obtener leyes universalmente válidas para todo tiempo y lugar, defendiendo que los fenómenos económicos forman parte de un momento histórico y, por tanto, la validez de las proposiciones se debe ceñir al entorno social en el que fueron observadas. Por lo tanto, defienden el método inductivo.
  \item \textit{Escuela marginalista.} Representada por autores como MENGER y WALRAS que apuestan por el método deductivo. A nivel metodológico también existen diferencias entre distintos colectivos marginalistas:
  \begin{la}
    \item MENGER y sus sucesores de la Escuela Austriaca optan por un proceso de deducción lógica y razonamiento verbal y no matemático; y, por otro lado,
    \item Otros autores marginalistas se alinearon más claramente con el razonamiento matemático (destaca el ejemplo de WALRAS y PARETO en la cátedra de Lausanne que siguen la senda de COURNOT\footnote{COURNOT fue profesor de economía política y matemáticas de AUGUSTE WALRAS, padre de LÉON WALRAS.}).
  \end{la}
\end{la}

\subsubsection{Positivismo lógica: Circulo de Viena (años 20's)}
La metodología de la ciencia entró en el siglo XX con el desarrollo del positivismo lógico, que dio al método científico unos fundamentos filosóficos. En esta corriente se unieron el razonamiento deductivo y el deseo positivista de dejar que los hechos hablaran por sí mismos. Tuvo su origen en un grupo llamado Círculo de Viena, e intentó formalizar los métodos de los científicos describiendo los que en realidad éstos seguían. 

Éstos sostenían que los científicos desarrollan una estructura deductiva (una teoría lógica) que lleva a formular proposiciones contrastables empíricamente. Sin embargo, una teoría deductiva sólo se considera verdadera una vez que se ha contrastado y verificado empíricamente. Por tanto, el papel del científico debe ser desarrollar estas teorías lógicas y contrastarlas.\footnote{Aunque existía un debate entre los positivistas lógicos sobre qué constituía la verdad, todos coincidían en que se descubriría por medio de la observación empírica.} 

El positivismo lógico imperó en la filosofía de la ciencia solamente durante la década de 1920, pero su influencia en la economía duró mucho más. En definitiva, esta corriente representó la culminación de la creencia de que el fin de la ciencia es establecer la verdad pero, desde entonces, la metodología de la ciencia se ha alejado cada vez más de esa idea. El primer alejamiento se debió a la preocupación por el aspecto de la verificación de la teoría positivista lógica

\subsubsection{Falsacionismo: POPPER} 
¿Puede una teoría probarse cierta para todas las situaciones en las que pueda concebirse su aplicación? ¿Es posible diseñar un experimento que pruebe el cumplimiento de una teoría de manera general, para todas las situaciones posibles? ¿Una teoría que no pueda ser probada falsa merece ser tomada en consideración? 

Estas preguntas definen el problema de la falsabilidad de una teoría. 

El filósofo vienés KARL POPPER (1902-1994) fue uno de los principales impulsores de la falsabilidad como requisito a cumplir por las teorías científicas. Decimos que una teoría es falsable cuando excluye al menos un estado de la naturaleza, de tal manera que sea posible verificar que tal estado de la naturaleza se ha producido o no. Y si efectivamente ese estado de la naturaleza se ha realizado, la teoría queda probada falsa. Así, aunque una teoría no pueda nunca probarse cierta en el sentido de comprobar que todos los estados de la naturaleza que predice se producen efectivamente, sí puede al menos probarse errónea. En economía y más generalmente en ciencias sociales, el problema de la falsabilidad se plantea a menudo. 

La Teoría de MARX sobre las fases de la historia (materialismo histórico\footnote{El materialismo histórico es una teoría de análisis propuesta por KARL MARX (1818-1883) y FRIEDRICH ENGELS (1820-1895) para comprender la historia humana desde el punto de vista de la lucha de clases sociales por el control de los medios de producción. Esta teoría proporciona una visión de la historia y la sociedad desde una perspectiva materialista, es decir, enfocándose en las condiciones materiales de la existencia humana, como la producción y las relaciones entre las clases sociales. Además, se opone a la interpretación tradicional burguesa, que comprende la historia como la historia de las ideas y de los grandes hombres.\\
MARX argumentaba que el motor de las sociedades humanas era la lucha entre diferentes clases sociales.\\
Para estos autores: Comunismo primitivo $\to$ Esclavismo $\to$ Feudalismo $\to$ Capitalismo $\to$ Comunismo}) es un ejemplo paradigmático de la crítica que POPPER hizo a las teorías no falsables. Si en un momento dado no se producían las transformaciones sociales previstas por esta teoría, sus defensores argumentarían que, aunque no se habían producido aún, se producirían más tarde o más temprano. Y si seguían sin producirse, se volvería a dar la misma respuesta.\footnote{He aquí la falsabilidad: tal teoría predice un estado de la naturaleza futuro en algún momento no especificado pero, como ese momento futuro no ha sido definido de manera clara, resulta imposible probarla falsa. Sin embargo, si esta teoría hubiese excluido una determinada situación o fenómeno futuro, en un momento determinado, sería al menos posible concebir que tal teoría fuese probada falsa (independientemente de que de hecho lo fuese o no).

No sólo la teoría marxista de la historia ha sido criticada en base a su falsabilidad: muchas otras teorías y modelos económicos han sido criticados por la ausencia de falsabilidad. Entre ellas, por ejemplo, encontramos la teoría cuantitativa del dinero, en la medida en que no define explícitamente el momento temporal en el que ha de producirse el aumento de la inflación, y, por tanto, no descarta la realización de un determinado estado de la naturaleza cuya observación habría de falsar la teoría (si aún no ha aumentado la inflación como consecuencia de un aumento de la oferta de dinero por encima del crecimiento del output, ya se producirá). Otro ejemplo de teoría objeto de la misma crítica es la teoría de las crisis financieras de MINSKY: \textit{si aún no se ha producido una crisis financiera como resultado de un periodo de estabilidad, se producirá en el futuro}.} Por esto, POPPER propone como buen método científico aquel que contempla la posible falsabilidad. FRIEDMAN y BLAUG son ejemplosclásicos de la aplicación de la metodología propuesta por POPPER a la ciencia económica, considerando que el propósito de ésta es abordar problemáticas concretas que puedan contrastarse empíricamente y realizar predicciones.

Cabe decir, que también se ha planteado el argumento contrario: no es necesario que un modelo sea falsable para que un modelo sea válido o útil, y de hecho se trata de una característica irrelevante. 

En cualquier caso, la hora de entender y aplicar un modelo económico, es preciso considerar la posibilidad de su falsabilidad como una característica deseable (como veremos más adelante) ya que sería estupendo que los problemas metodológicos pudieran resolverse de una manera tan clara como sugiere el enfoque de POPPER, pero los debates metodológicos no son tan claros.
De hecho, no debemos perder de vista el hecho de que modelos no falsables también son o han sido capaces de predecir y explicar fenómenos económicos, y que la necesidad de falsabilidad es en todo caso objeto de debate:\footnote{De hecho, las tendencias más recientes han alejado la metodología cada vez más de esas nítidas distinciones.}
\begin{lnum}
  \item En primer lugar, las predicciones empíricas de algunas teorías no pueden contrastarse porque no existe la tecnología necesaria para contrastarlas, ¿Qué debe hacerse con esas teorías?
  \item En segundo lugar, es difícil saber cuándo se ha falsado o no una teoría;\footnote{Por ejemplo, si una contrastación empírica no produce los resultados esperados, el investigador puede atribuir, y a menudo atribuye, el fracaso a fallos del procedimiento de contrastación o a algún factor exógeno. Por tanto, una sola contrastación empírica negativa a menudo no invalida la teoría}
  \item El tercer problema se debe a la manera de pensar de los investigadores, que pueden no contrastar las implicaciones de una teoría establecida, suponiendo que son verdaderas. Ese modo de pensar puede llevar a no aceptar teorías nuevas y posiblemente más defendibles; y,
  \item Un cuarto problema, será la crítica de DURHEM, quien argumentará que es imposible falsar de forma concluyente una hipótesis científica concreta porque a la hora de enfrentarla a la evidencia empírica siempre estaremos contrastando la hipótesis científica concreta con un conjunto de hipótesis auxiliares que la acompañan implícitamente, de manera que nunca podremos estar seguros de si la refutación se debe a problemas en la hipótesis principal o en las auxiliares. Esto supone un problema operativo que POPPER admite, dado que los científicos se verán tentados a revisar de forma continuada las hipótesis auxiliares con el objetivo de defender sus teorías frente a los intentos de refutación de las mismas. Ante esta situación, POPPER propondrá completar su criterio de demarcación de modo que lo que permitirá distinguir las teorías científicas de las no científicas, según él, será su posibilidad de falsación más la imposibilidad de establecer \textit{supuestos auxiliares ad hoc} (a los que POPPER denominará \textit{estrategias inmunizadoras}).
\end{lnum}

\subsubsection{Paradigmas científicos: \authorfont{KUHN}}
¿Cómo cambian la forma de generar nuevo conocimiento? ¿Los avances científicos aparecen a una velocidad lineal y de manera acumulativa, o el nuevo conocimiento científico se concentra en periodos de rápido avance y después tienen lugar periodos de estancamiento? ¿Cómo caracterizar ese proceso? 

El filósofo estadounidense THOMAS KUHN (1922-1996) propuso el concepto de paradigma científico como el conjunto de supuestos esenciales, métodos habitualmente aplicados y conjunto de principios éticos, científicos y experimentales que están presentes en la ciencia practicada en un momento determinado y en una comunidad dada.\footnote{Es necesario tener presente que el concepto de paradigma científico es en cualquier caso una herramienta historiográfica que tiene por objetivo entender la evolución del conocimiento científico, la relación entre modelos y teorías y la persistencia de determinadas teorías a pesar de su falsación total o parcial (si es que es posible). Así, el concepto de los paradigmas científicos no plantea éstos como una proposición a priori a la generación de nuevo conocimiento científico.\\
Así, sería erróneo partir del concepto de paradigma e imaginar que todas las comunidades científicas siguen los principios y elementos del paradigma dominante como una acción consciente y deliberada, y que cuando ocurre una revolución científica ello ocurre porque un grupo de científicos han tomado la decisión de revolucionar el paradigma dominante. Más bien, los paradigmas deben aplicarse a posteriori para entender el contexto en el que un determinado corpus de conocimiento científico se desarrolló, y cuáles fueron las causas, en ese contexto, de la aparición de un nuevo paradigma con teorías y modelos que explicasen mejor determinados fenómenos.}

\imagenfit{FasesKUHN.png}{Fases de un paradigma científico. KUHN, 1962}{\href{https://es.wikipedia.org/wiki/La_estructura_de_las_revoluciones_científicas}{la estructura de las revoluciones científicas. Wikipedia}}

El nuevo conocimiento adquirido en el contexto de ese paradigma es definido por el autor como \textit{ciencia normal}. Sin embargo, en todo paradigma científico existen algunos fenómenos que no han sido explicados de manera satisfactoria (o que no han sido explicados en modo alguno con las herramientas de ese paradigma). Cuando estos fenómenos se generalizan o se vuelven lo suficientemente importantes como para poner en duda los métodos, supuestos y teorías del paradigma en cuestión, puede aparecer un nuevo paradigma dominante en lo que KUHN denominó revolución científica.\footnote{En economía encontramos numerosos ejemplos de estas revoluciones. La aparición de la economía neoclásica, que aplicaba herramientas tales como el cálculo diferencial o supuestos como la racionalidad optimizadora del bienestar supuso una revolución frente a los paradigmas previamente dominantes en Europa que primaban las teorías esencialmente formuladas en términos narrativos y que concebían a menudo los fenómenos económicos como el resultado de leyes más o menos mecánicas que encuadraban el comportamiento de los seres humanos que tomaban decisiones con contenido económico.}

Por tanto, mientras que, para POPPER, la verdad (o lo más que nos podemos acercar a ella) vencerá y por lo tanto estamos en un estado de revolución permanente, para KUHN, es posible que exista una teoría superior, pero no se adopta debido a la inercia que favorece al paradigma existente. Por tanto, la teoría reinante no es necesariamente la mejor.\footnote{Para ver algunos ejemplos de aplicación en Economía de la Metodología o Enfoque propuesto por KUHN se puede ver en [\authorfont{Ver Anexo \ref{anx:paradigmaKUHN}}]}

\subsubsection{Programas de investigación: \authorfont{LAKATOS}}  
Si es necesario que las teorías científicas sean falsables, y ocurre de hecho que son falsadas, ¿es preciso desechar todo el conjunto de teorías relacionadas con aquélla que es falsada? En la práctica, ¿desechamos los paradigmas que contienen una o varias teorías que son probadas falsas?

IMRE LAKATOS (1922- 1974) planteó el concepto de programa de investigación como un intento de compatibilizar el requisito de falsabilidad y el concepto de paradigma científico que perdura en el tiempo.\footnote{Si bien el argumento a favor de la falsabilidad planteado por POPPER se estableció rápidamente como un requisito deseable de las teorías científicas y en cierto grado, también en economía, el concepto de paradigma de KUHN alcanzó también notable aceptación. En este contexto, y tomando en consideración la evidente persistencia de paradigmas que no son capaces de explicar sin lugar a dudas todos los fenómenos observados o que contienen teorías que son probádas falsas una o varias veces (algo de lo que los economistas constatamos con notable frecuencia), se manifiesta una aparente contradicción entre lo deseable de desechar teorías falsadas y lo deseable de no desechar inmediatamente y casi de manera constante los paradigmas científicos prevalentes en un momento determinado.} 

Para LAKATOS, un programa de investigación es la combinación de un núcleo científico compuesto por una serie de teorías no falsables que no pueden ser abandonadas sin abandonar por completo el programa de investigación, y un cinturón protector en el que se encuadran un conjunto más o menos amplio de teorías, corolarios, prescripciones, experimentos, etc. que sí pueden ser falsados. Cuando uno de esos elementos del cinturón protector resulta falsado, se su puede sustituir por otro que explique mejor un determinado fenómeno, sin hacer necesario desechar el núcleo del programa de investigación. Así, el cinturón protector sirve de escudo ante las anomalías que escapan en un primer momento a las explicaciones propuestas para entender un determinado fenómeno. 

Las teorías económicas son a menudo organizadas como programas de investigación. Tomemos, por ejemplo, el programa de investigación de la llamada Nueva Macroeconomía Clásica. En este programa, el núcleo está formado por supuestos tales como la optimización racional basada en la Hipótesis de Expectativas Racionales, herramientas matemáticas tales como la programación dinámica o los métodos recursivos de optimización, y la crítica de LUCAS como razón de ser del propio programa. En el cinturón protector, encontramos modelos como el modelo de relación entre producción, precios y oferta monetaria de LUCAS (1972), el modelo de inconsistencia de la política económica en un contexto dinámico de KYDLAND y PRESCOTT (1977) o el Modelo del Ciclo Real en sus múltiples variantes. Estos elementos del cinturón protector pueden ser probados falsos empíricamente, pero ello no obsta para que los elementos del núcleo sean utilizados para formular nuevas teorías que utilicen las mismas herramientas metodológicas y que tengan similares objetivos esenciales.

\subsubsection{Anarquismo metodológico: \authorfont{FEYERBAND}}
El anarquismo metodológico o epistemológico de PAUL FAYERABEND (1924-1994) sostiene que no existen principios metodológicos de aplicación general y sin excepciones que gobiernen el proceso de avance del conocimiento que constituye el objeto de la ciencia. 

Concebir el método científico como un conjunto de reglas fijas y universales acaba siendo perjudicial para el avance del conocimiento científico porque tales reglas acaban por limitar la actividad científica y, por ende, el progreso del conocimiento científico. Tampoco es deseable la imposición de reglas rígidas que constriñan el comportamiento de los científicos. Así, lo óptimo es admitir un cierto grado de anarquismo metodológico que deje margen a la comunidad científica para explorar nuevas herramientas y procesos de asimilación de nuevo conocimiento que hagan posible el avance de la ciencia en el marco de las teorías ya existentes, y la designación de nuevos objetivos científicos.

\subsubsection{Controversia de los supuestos} 
¿Deben los modelos económicos basar sus implicaciones en supuestos realistas? ¿El grado de realismo de los supuestos determina la bondad de un modelo a la hora de comprender y predecir un fenómeno económico? ¿Adoptar supuestos poco realistas es aceptable si ello permite formular modelos precisos y tratables? El debate en torno a estas preguntas ha sido una constante en la ciencia económica desde el momento en que se generaliza la formulación de modelos. En el siglo XIX, Stuart Mill, Cairnes y Nassau Sr. mantuvieron una intensa querella en torno a esta cuestión. 

Cairnes afirmaba que una teoría económica era cierta en la media en que estuviese basada en premisas cuya verdad fuese innegable. En la medida en que las premisas dejasen de cumplir tales condiciones, habría que ajustarlas para que se adaptasen lo más posible a la realidad y la teoría pudiese así serlo también. Esta postura de Cairnes prevaleció más o menos indiscutida hasta mediados del siglo XX. 

En 1953, Milton Friedman publica un ensayo5 en el que defiende una tesis incompatible a la hegemónica en aquel momento. Friedman afirmaba que un modelo económico no es menos válido porque no se conozca la verdad de todos los supuestos que subyacen a las predicciones. Así, no es necesario que todos los supuestos de un modelo sean verdaderos, y no hay que basar el juicio de validez de un modelo en la veracidad de los supuestos de los que se parte.\footnote{Para que un modelo sea útil, basta con que algunos de los supuestos sean ciertos, y que la secuencia lógica que lleva de los supuestos a las conclusiones sea formalmente correcta y materialmente capaz de predecir un resultado que efectivamente tiene lugar. La tesis de Friedman fue criticada por autores como Koopmans o Samuelson, entre otros. Koopmans interpretó la postura de Friedman como \textit{el realismo de los supuestos de un modelo es totalmente irrelevante, y lo único importante es la consistencia de la predicción con lo observado empíricamente}. Esta reinterpretación fue a su vez objeto de críticas por distorsionar notablemente la tesis postulada por Friedman.} 

En la práctica, la controversia de los supuestos es probablemente menos relevante que lo que se desprende de la intensidad que alcanzó el debate. Los modelos económicos han seguido y siguen basándose parcialmente en supuestos falsos, de la misma forma que sucede en el resto de las ciencias sociales. La complejidad de los fenómenos humanos es inasequible al reducido número de supuestos que pueden abarcarse en un modelo económico tratable. La simple omisión del elevadísimo número de circunstancias que definen a un ser humano que interacciona en un contexto económico constituye en sí misma un supuesto irrealista que resulta inevitable adoptar: ¿cuántas formas de capital físico o financiero existen?, ¿tienen los agentes optimizadores del ocio y el consumo animales domésticos que afecten al resultado de ese cálculo optimizador?


\subsubsection{Otros desarrollos}

\paragraph{Enfoque sociológico}
El enfoque sociológico de la ciencia caracteriza el éxito o fracaso de una teoría científica como el resultado de la confluencia de los intereses de diferentes grupos sociales a los que beneficia o perjudica el éxito o el fracaso de una teoría u otra. Así, este enfoque no trata de entender el valor de una teoría en términos de su capacidad para explicar o predecir un fenómeno. Más bien, se trata de explicar la prevalencia de una teoría respondiendo a preguntas tales como ¿a quién interesa que esta teoría se imponga en un determinado contexto social, académico, social, etc.?, ¿qué incentivos a investigar dentro de un determinado paradigma tienen los miembros de una comunidad científica?, o ¿qué incentivos tienen los miembros de esa comunidad académica para salirse del paradigma predominante, desarrollando nuevas teorías dentro de paradigmas distintos? Así, los incentivos económicos son un punto esencial en este análisis del método científico: ¿qué condiciones de financiación, influencia política y académica, reconocimiento social, etc. pueden obtener los científicos que investiguen en una dirección u otra? 

Este enfoque tiene uno de sus orígenes más inmediatos en la obra de KARL MARX, pero ha tenido un muy variado desarrollo hasta la actualidad. Dentro de la propia ciencia económica, autores como MILTON FRIEDMAN o corrientes como la escuela del Public Choice, entre otras, han planteado las luchas por la financiación para investigación, las posiciones dentro de departamentos académicos o la influencia en los programas de gasto de la administración pública, como factores que determinan que una u otra teoría acabe imponiéndose en términos de su aceptación generalizada por la comunidad académica en un lugar determinado.


\subsection{Métodología convencional: Principales elementos}

\subsubsection{Supuestos y axiomas}


\paragraph{Análisis marginal}

\paragraph{Escasez}


\paragraph{Individualismo metodológico}
¿Los fenómenos económicos que observamos a nivel agregado son el resultado directo de las decisiones individuales de los agentes económicos? O, por el contrario, ¿existen fenómenos cuya comprensión a partir de la decisión individual resulta imposible, y sólo resulta posible explicarlos como resultado de dinámicas que ocurren a nivel agregado? Las diferentes respuestas a estas preguntas dan forma a un debate distinto, pero tangencialmente conectado con la crítica de Lucas y la necesidad de microfundamentar los modelos microeconómicos:
\begin{la}
    \item El individualismo metodológico es la doctrina según la cual la explicación dada sobre fenómenos sociales se basa en el comportamiento de agentes individuales; y,
    \item Por otro lado, el holismo metodológico afirma que existen leyes naturales de nivel macroscópico que se aplican al conjunto de la sociedad y que explican fenómenos que trascienden a la decisión individual y son imposibles de explicar a partir de procesos de decisión individual.
\end{la}

La hegemonía en el contexto académico de la modelización microeconómica y a mayores, la microfundamentación de los modelos macroeconómicos, ha supuesto en las últimas décadas una victoria del individualismo metodológico en la economía.\footnote{Esta dicotomía metodológica no es exclusiva de la economía, sino que es común en la mayoria de las ciencias sociales (en especial, en sociología).} Sin embargo, el fenómeno de la emergencia\footnote{Es decir, la aparición de fenómenos que sólo pueden ser descritos y entendidos desde el punto de vista agregado.} sigue presente tanto en los fenómenos observados como en las explicaciones propuestas para estos. 

Los modelos que incorporan dinámicas caóticas o los modelos basados en agentes -más relevantes en la actualidad- son ejemplos de que la existencia de la emergencia y la necesidad de  explicarla siguen siendo una aspiración de la economía, y el individualismo metodológico no es suficiente para explicarlo.

\subsubsection{Instrumentos analíticos}

\paragraph{Optimización: Matemáticas}

\paragraph{Estadística}


\paragraph{Econometría}


\subsubsection{Cualidades deseables de un modelo} 
Como hemos analizado, el debate sobre qué características concretas deben satisfacer los modelos económicos trasciende las épocas y las disciplinas dentro de la economía.\footnote{En definitiva, un modelo económico debe servir para comprender el porqué de un fenómeno, y para realizar predicciones sobre futuras manifestaciones de éste. Alcanzar este objetivo de manera satisfactoria implica cumplir múltiples requisitos, que a veces se solapan y en la mayoría de los casos no pueden cumplirse todos a la vez.} Sin embargo, algunas de ellas pueden aceptarse con cierta pretensión de generalidad. GABAIX y LAIBSON (2008) proponen siete cualidades cuya presencia es deseable en cualquier modelo económico:\footnote{No son cualidades suficientes ni absolutamente necesarias para que un modelo deba utilizarse para explicar un fenómeno, ni en la práctica es posible formular modelos que cumplan todas ellas, pues aparecen habitualmente trade-offs difíciles de resolver entre dos o más cualidades. 

Sin embargo, esta lista de cualidades es una aproximación razonable generalmente a los problemas que enfrenta cualquier economista que trate de comprender un fenómeno económico mediante la formulación de un modelo, y debe resolver preguntas tales como: ¿la explicación de mi modelo es demasiado complicada para un fenómeno cuya explicación es realmente trivial?, ¿mi modelo puede probarse falso, o siempre se puede argumentar que será posible probarlo cierto en el futuro si no ha sido posible hacerlo aún? ¿las conclusiones de mi modelo pueden generalizarse a un número elevado de situaciones, aunque no sean totalmente idénticas?}
\begin{lnum}
    \item \textit{Profundidad conceptual.} Un modelo económico debe explicar fenómenos cuyas causas y consecuencias no estén trivialmente relacionadas. Si un fenómeno tiene una explicación evidente, no son necesarias explicaciones complejas o excesivamente formalizadas: el coste en términos intelectuales y de tiempo es simplemente demasiado elevado para entender un fenómeno cuyas causas saltan a la vista para cualquier observador.\footnote{En las últimas décadas, la ausencia de profundidad conceptual en relación con la complejidad técnica de algunos modelos económicos ha sido objeto de crítica creciente. Estas críticas se han extendido también al uso de técnicas matemáticas complejas como barrera de entrada para el acceso a la investigación económica y al policy-making.}
    \item \textit{Falsabilidad}. Como se explicó en el segundo apartado de este tema, un modelo es falsable cuando excluye algún estado de la naturaleza en el futuro: si tal estado de la naturaleza llega a materializarse, ese modelo podrá considerarse falso.\footnote{Si el modelo no excluye ningún estado futuro, es imposible considerarlo como falso. Supongamos, por ejemplo, un modelo que predijese que la actividad económica va a colapsar en el futuro como resultado de determinadas políticas económicas. Y supongamos que ese \textit{el futuro} se define de manera lo suficientemente vaga como para que el modelo (o los defensores de la capacidad del modelo como herramienta predictiva) siempre pueda contestar a la pregunta de \textit{¿cuándo colapsará la actividad económica?} con un \textit{Aún no, pero lo hará sin duda en el futuro}. En este, caso estamos ante un modelo no falsable, que nunca puede probarse cómo falso. Si el estado de la naturaleza predicho por el modelo se materializa finalmente, el modelo habrá acertado. Pero si no se materializa será porque \textit{hay que esperar un poco más...}; Un modelo con estas características, ¿sería útil?}
    \item \textit{Generelizabilidad.} Un modelo es generalizable cuando es capaz de explicar un rango amplio de fenómenos pasados a partir de modificaciones de las condiciones iniciales del modelo que no requieren cambios en su estructura interna.\footnote{En la práctica, la generalizabilidad de un modelo determina en gran medida la difusión del modelo como herramienta explicativa. 
    
    Tomemos como ejemplo el conocido modelo de MUNDELL-FLEMING. Con apenas un cambio en los supuestos relativos a la variable que no cambia en una de las tres ecuaciones, es posible formular predicciones para todas las combinaciones de economías con tipo de cambio fijo o flexible, y movilidad de capitales o ausencia de ella. Si economías con tipo de cambio flexible requiriesen un modelo totalmente distinto al necesario para representar economías con tipo de cambio fijo, estaríamos ante un modelo mucho menos generalizable. Casi todos los modelos económicos más conocidos son generalizables a un número relativamente elevado de situaciones.
    }
    \item \textit{Tratabilidad}. Un modelo es tratable cuando la predicción de estados de la naturaleza futuros puede llevarse a cabo aplicando métodos al alcance de cualquier persona con un nivel medio de conocimiento en la materia objeto de predicción.\footnote{La tratabilidad de un modelo es también una condición subjetiva: un modelo económico que requiera del uso de técnicas tales como el control óptimo o la topología diferencial, no es tratable si sus usuarios son estudiantes de bachillerato en su primer contacto con la economía. Pero puede serlo para estudiantes de doctorado. 
    
    En la práctica, se trata de una característica esencial para que un modelo sea aplicado con solvencia: si quien debe formular predicciones o entender fenómenos no es capaz de entender la lógica interna del modelo para derivar una conclusión de unas condiciones iniciales o unas premisas, se estará limitando a aplicar una regla arbitraria de correlación de fenómenos. En la actualidad, el desarrollo de la computación y el acceso generalizado a herramientas de software que permiten resolver modelos económicos ha difuminado en cierta medida la importancia de la tratabilidad de un modelo. Pero si quien utiliza un modelo no puede entender su funcionamiento interno por no ser tratable para este usuario, tampoco será capaz de criticarlo o encontrar errores en el modelo, ni compararlo con otros modelos alternativos.
    }
    \item \textit{Parsimonia}. Un modelo es parsimonioso cuando utiliza la mínima cantidad de información necesaria para explicar un fenómeno concreto.\footnote{Así, un modelo parsimonioso se limita a tomar los mínimos supuestos necesarios para derivar unas conclusiones que expliquen el fenómeno observado, evitando incluir información superflua. El cuento \textit{Del rigor en la ciencia} de BORGES plantea una parábola que ilustra la necesidad de formular modelos parsimoniosos (\textit{Sylvie and Bruno} de CARROLL contiene una parábola). 
    
    En un imperio antiguo, los cartógrafos del emperador construían mapas del imperio tan exactos, que su superficie era tan grande como el propio imperio. Con el tiempo, las siguientes generaciones se dieron cuenta de la inutilidad de un mapa de esas características: un mapa parsimonioso, con mucha menos información, habría sido mucho más útil a la hora de guiar los desplazamientos de los habitantes del imperio. 
    
    Sin embargo, y a pesar de lo deseable de la parsimonia de un modelo, existen límites a esta característica: más allá de cierto límite, menos información en un modelo resulta en conclusiones diferentes, o en el incumplimiento de otras características deseables. En economía, la parsimonia es una de las principales razones para utilizar lenguaje formal (como el lenguaje matemático) en la modelización. La representación formal de un modelo permite definir de manera clara qué es lo que realmente se está suponiendo para construir el modelo, y excluir todo aquello que no se menciona expresamente.
    }
    \item \textit{Consistencia empírica}. Un modelo es empíricamente consistente cuando sus predicciones son consistentes con los fenómenos efectivamente observados.\footnote{Si un modelo predice que tales condiciones iniciales inducirán tal resultado, y ese resultado efectivamente tiene lugar, estaremos ante un modelo empíricamente consistente. En el contexto económico podemos imaginar un modelo de inflación de la siguiente forma: ante una tasa de desempleo por debajo de una tendencia estimada con observaciones pasadas, la inflación tiende a acelerarse.
    
    A primera vista puede entenderse esta cualidad como una condición \textit{sine qua non} para utilizar un modelo económico, en la práctica no siempre es sencillo ni apropiado desechar un modelo porque no haya sido capaz de predecir.
    }
    \item \textit{Precisión predictiva}. Un modelo predice con precisión cuando sus estimaciones de futuros estados de la naturaleza discrepan poco o nada, en términos cuantitativos, con lo que efectivamente sucede.\footnote{En la práctica, es imposible postular esta cualidad para muchos modelos económicos, por el simple motivo de que no formulan ninguna predicción en términos cuantitativos precisos, limitándose a predecir o explicar un resultado cualitativo. Por ejemplo, una apreciación o depreciación de la moneda, una caída de la producción o un aumento de la inflación. El desarrollo de la estadística, la econometría, los métodos de recogida de datos y el aumento de la capacidad de los microprocesadores de los últimos setenta años ha permitido el desarrollo de modelos que tratan de estimar de manera precisa estados futuros de las economías dadas unas condiciones iniciales, y unas políticas públicas o unos cambios determinados en variables demográficas, tecnológicas, políticas o sociales, entre otras.
    }
\end{lnum}

\subsection{Debates y Conceptos de largo alcance}


\subsubsection{La existencia del Equilibrio}




\subsubsection{Coste de oportunidad: Optimización}
\textcolor{red}{\textbf{OJO} Este concepto es fundamental, hay que explicarlo bien. Es, en puridad, la gran cuestión de la economía como ciencia normativa y la que causa más estragos en la economía como ciencia descriptiva}

La relación entre economía y psicología no es en absoluto reciente. En \textit{La Teoría de los Sentimientos Morales} de ADAM SMITH (1759), se analizan las causas psicológicas de numerosas instancias de intercambio económico. En \textit{La Riqueza de las Naciones} (1776) se proponen también explicaciones de corte psicológico para variados fenómenos económicos como la especialización en el trabajo, la determinación del nivel salarial, la tributación o la innovación empresarial. 

Otros autores del periodo clásico tales como MALTHUS, BENTHAM o MILL trataron de describir comportamientos humanos que, si bien no podían desviarse de la optimización racional por aún no haberse desarrollado este concepto, sí que pueden enmarcarse en el uso de la introspección psicológica que constituye uno de los elementos centrales de la economía conductual. Es, sin embargo, en el siglo XX cuando aparece la economía conductual con este nombre. 

En los años 50, HERBERT SIMON (que recibiría el Premio Nobel en 1978) inicia un programa de investigación pionero que tiene por objetivo describir las limitaciones cognitivas del ser humano en la toma de decisiones económicas como resultado del elevado coste computacional que tales cálculos tienen para el cerebro. Aunque se postula que en general los agentes económicos sí son racionales en el sentido de que actúan de manera consistente con la consecución de sus deseos, no siempre resulta posible determinar cuál es de hecho la decisión óptima por el prohibitivo coste que pueden suponer dos factores: (i) el proceso de introspección necesario para conocer las propias preferencias; (ii) las alternativas de decisión; y, (iii) las probabilidades que cada una de esas decisiones asignan a los distintos estados de la naturaleza objeto de las preferencias de los agentes.

Este programa de investigación, junto con los trabajos de otros autores tales como ALLAIS, ELLSBERG o FELLNER, tuvieron una influencia esencial en el trabajo seminal de KAHNERNAN y TVERSKY en los años 70's que formalizaría lo que se conoce como Economía Conductual.


\subsubsection{Mercados: Última salva de la racionalidad limitada}
\textcolor{red}{\textbf{OJO} Este elemento es esencial, hay que explicarlo muy bien. Para una introducción de parte se puede recurrir a THALER (2021). En eencia, es el \textit{último bastión} en el que se refugia la ortodoxia económica para refugiarse o conciliar el tsunami de la economía conductual, primera corriente heterodoxo que ha logrado poner entre la espada y la pared a la corriente dominante

\textbf{LA DISCIPLINA DEL MERCADO:} Independientemente de cómo actúen \textit{algunos agentes}, el mercado disciplina su comportamiento y verifica el cumplimiento de los supuestos de partida de una forma agregada. ¿A sí? -> clé

}






\clearpage
\section{¿Ciencia Económica?: Cuestiones y debates actuales}
\puntosclave{
    La economía continúa avanzando y madurando como ciencia autónoma, generando un cuerpo de conocimiento dinámico.
    
    Los principales avances son en materia de método.
}

\subsection{La Economía Conductual} 
La economía conductual es un conjunto de herramientas metodológicas que tienen por objeto explicar y predecir decisiones humanas con contenido económico en las que los individuos se desvían de una noción de racionalidad,\footnote{Entendiendo racionalidad como la toma de decisiones e implementación de las mismas de manera tal que el agente decisor alcance o al menos se acerque a la consecución de sus preferencias, dadas éstas y un conjunto de infonnación a su alcance.} y que han sido constatadas empíricamente de manera experimental u observacional. Con el objetivo de desarrollar tales herramientas metodológicas, la economía conductual ha incorporado numerosas técnicas y conceptos provenientes de la psicología, así como adoptado un enfoque eminentemente experimental a la hora de obtener nuevo conocimiento.

\subsubsection{Aportaciones notables: regularidad empírica}
El producto principal de la economía conductual es un compendio de estilizaciones de desviaciones sistemáticas de la racionalidad, para las que se postula un modelo subyacente que lo explica. Sin pretensión de exhaustividad, resulta ilustrativo presentar algunas de las más conocidas desviaciones de la racionalidad que han sido descritas en la literatura: 
\begin{lnum}
    \item \textit{Framing effect.} La decisión entre diferentes alternativas está condicionada por la presentación que se hace de éstas al agente decisor. Esta presentación puede modularse de distintas formas: verbales, visuales, temporales, sensoriales. . . Así, puede mostrarse empíricamente como ante alternativas de decisión materialmente idénticas. los seres humanos eligen unas u otras en función de cómo hayan sido presentadas. En el impacto que tenga unapresentación u otra influyen factores innatos a la capacidad cognitiva del agente decisor. pero también factores culturales y sociales.
    \item \textit{Preferencias no lineales en la probabilidad.} En el contexto de la decisión en presencia de riesgos objetivamente cuantificables, los seres humanos no ponderan en igual medida idénticas variaciones relativas de la probabilidad de que ocurra un determinado suceso. cuando el nivel absoluto de probabilidad cambia. Por ejemplo, empíricamente puede mostrarse que un aumento del 1 \% de la probabilidad de que ocurra un suceso catastrófico se pondera en mayor medida (es decir, se exigen mayores pagos o compensaciones en contrapartida) cuando la probabilidad inicial de que ocurriese tal suceso era 0\%, que cuando era ya del 40 o el 50\%. ¿Tiene el mismoimpacto un aumento del riesgo de muerte del 1\% cuando el riesgo es inicialmente 0\%, o cuando es inicialmente 57\%? ¿Estaríamos dispuesto a aceptar la misma compensación por pasar del 0\% al 1\%, que del 56\% al 57\%?
    \item \textit{Descuento hiperbólico.} Los seres humanos deben decidir no sólo entre alternativas de realización inmediata, sino entre opciones que inducen bienestar y malestar en diferentes momentos temporales futuros.\footnote{Así, la decisión en contexto intertemporal implica una dimensión adicional representada por el propio tiempo: ¿es lo mismo obtener un pago con un valor real de 10.000€ dentro de un año o dentro de tres años? Para representar esta dimensión temporal, es habitual aplicar un factor de descuento decreciente con el tiempo, de tal manera que cuánto más alejado en el tiempo esté un suceso, menor sea su peso en la decisión del agente. La senda concreta que sigue este factor de descuento depende de la importancia relativa del presente respecto al futuro, (capturada habitualmente por un parámetro denominado tasa de descuento subjetivo), y por las propiedades que se quiera atribuir a las preferencias en la decisión en contexto intertemporal. Las propiedades de estacionariedad (la decisión entre dos alternativas a no depende del periodo concreto, sino de la diferencia temporal entre periodos y entre cuantías obtenidas) y de consistencia temporal (la decisión que era óptima en el periodo $t$ sigue siendo óptima en $t+n$ si la información disponible es idéntica) se consideran habitualmente como propiedades deseables para representar un proceso de decisión coherente y predecible. El llamado descuento exponencial, en el que la tasa de descuento equivale a $\frac{1}{(1+\rho)^t}$, satisface las propiedades de estacionariedad y consistencia temporal. Sin embargo, existe evidencia empírica que muestra que la decisión humana en contexto intertemporal puede aproximarse mejor con un descuento hiperbólico, en el que la tasa de descuento equivale a $\frac{1}{(1+\rho \; t)}$. En esta forma de descuento, las propiedades de consistencia temporal y estacionariedad ya no se cumplen. Un ejemplo que ilustra las diferentes propiedades parte de la siguiente decisión: obtener cincuenta euros ahora o bien obtener cien dentro de un año. Ante tal disyuntiva, empíricamente una mayoría prefiere los cincuenta euros ahora. Sin embargo, ante otras dos alternativas tales que con una se obtienen cincuenta euros dentro de cinco años o cien euros dentro de seis años, una mayoría prefiere los cien euros dentro de seis años. Si las preferencias cumpliesen la propiedad de estacionariedad, la decisión debería ser igual en ambos horizontes temporales (ahora y dentro de un año, dentro de cinco y dentro de seis años), porque lo único que ha cambiado son los periodos concretos, pero no la diferencia entre periodos. Así, un descuento exponencial no podría dar lugar a tal observación empírica, mientras que un descuento hiperbólico sí podría generar tal resultado empírico.}
    \item \textit{Dependencia del origen de la incertidumbre.} La decisión humana tiene en cuenta el origen de la incertidumbre cuando debe elegir entre dos alternativas inciertas. Y tiende a preferir situaciones en las que el riesgo es cuantificable y conocido, frente a situaciones en las que la distribución subyacente se desconoce totalmente.\footnote{La Paradoja de ELLSBERG es un ejemplo paradigmático de este fenómeno. En esta paradoja, se parte de dos bolsas con bolas de dos colores. En la primera de las bolsas, la distribución relativa de los colores de las bolas es conocida. En la segunda bolsa, la distribución relativa de colores se desconoce totalmente. Si se presenta a un grupo de individuos con la posibilidad de ganar una cantidad dada si sacan una bola de un color, y se les permite elegir la bolsa de donde ha de salir la bola, la mayoría prefiere la primera bolsa (la que tiene una distribución conocida). A continuación, se les presenta con otra apuesta, en la que se les paga la misma cantidad si sacan una bola de un color diferente al de la primera ocasión y se les permite también elegir la bolsa de donde sacar las bolas. En este segundo caso, una mayoría vuelve a preferir la primera bolsa. Ambas decisiones en conjunto implican una preferencia sistemática por apuestas con distribuciones de probabilidad conocidas, dado que no existe una distribución de probabilidades objetivas implícita que pudiese racionalizar tal decisión entre las dos apuestas.}
    \item \textit{Aversión a la pérdida.} La decisión humana muestra una tendencia a sobreponderar la importancia de las pérdidas respecto a ganancias de igual cuantía.\footnote{El rango de ganancias y pérdidas para el que se manifiesta el fenómeno, así como la importancia relativa de tales pérdidas y ganancias respecto a la riqueza previa, muestra variación en función del contexto y los agentes. La aversión al arrepentimiento es otro fenómeno relacionado con la aversión a la pérdida. En este caso, los agentes valoran negativamente la posibilidad de saber a posteriori que, si hubiesen tornado otra decisión, su utilidad habría sido mayor. Y por ello, adoptan la decisión que minimiza el posible arrepentimiento, aunque en términos de utilidad esperada no se tratase necesariamente de la decisión óptima.}
\end{lnum} 

\subsection{La Economía de la Complejidad}



\subsection{La Economía Experimental}
El desarrollo de la economía experimental, así como el acceso a conjuntos de microdatos y la creciente capacidad de procesamiento de éstos ha inducido la consolidación de la economía conductual como un enfoque transversal a prácticamente todos los fenómenos microeconórnicos esenciales: 1a decisión entre alternativas, las decisiones de consumo, la organización de la empresa y sus decisiones de producción, la distribución de recursos escasos en un  contexto intertertemporal, la decisión en contexto de incertidumbre, o la toma de decisiones en presencia de interdependencia estratégica.

\subsection{Microdatos}

\subsubsection{Heterogeneidad entre agentes económicos}

\subsubsection{Desigualdad}


\subsubsection{Colas de la distribución}



\clearpage
\section{Conclusión}

\subsection{Recapitulación}



\subsection{Extensiones y relación con otras partes del Temario}



\subsection{Opinión}

\subsubsection{Idea final (Salida o cierra)}

\clearpage
\pagenumbering{gobble}
\MyBibliography
\MyBibItem{Autor}{Título}{Año}{DOI -- LINK}


% --- Anexos (no aparecen en el índice) ---
\clearpage
\anexo{\textbf{Preguntas Test}}
%\input{\TemaCode/questions.tex} 





\clearpage
\anexo{Evolución histórica del objeto de la Economía}\label{anx:historiaobjeto}

El análisis de esta cuestión desde una perspectiva histórica permite poner en contexto la evolución de los objetos de la economía con otras disciplinas, así como con el contexto social y político.

\textsc{Economía como análisis moral del intercambio}

La dimensión moral y ética de los intercambios económicos ha estado presente con mayor o menor importancia en toda la historia del pensamiento. En la Europa medieval, el pensamiento escolástico centró su análisis en los aspectos morales del intercambio, tratando de definir con claridad qué transacciones eran o no compatibles con el derecho natural. Destaca la obra de TOMÁS DE AQUINO y en concreto, sus análisis relativos al precio justo, al uso de la propiedad privada, al crédito y la usura, y a la valoración ética de las diferentes fonnas de trabajo. En el mundo islámico, es especialmente relevante la obra de IBN KHALDUN (s. XIV) por su análisis pionero en cuestiones tales como la teoría del valor (que precede a SMITH y a RICARDO y se inspira en la teoría del intercambio justo de ARISTÓTELES), economía monetaria y crecimiento económico en el contexto de su compatibilidad con la religión islámica.

\textsc{Economía como ciencia de buen gobierno}

KAUTILA en la India clásica, TÁCITO en el siglo I d. C, o CONSTANTINO VII en el Imperio Bizantino, -entre otros-, compilaron obras en las que se aconsejaba al gobernante cómo conducir su actuación para fortalecer el poder del estado frente a otros estados, sus propios súbditos, y otros aspirantes a situarse a la cabeza del poder estatal. Esto es, cómo alcanzar el buen gobierno. Y alcanzar tal objetivo implicaba necesariamente implementar políticas de contenido económico.

En el contexto del surgimiento de organizaciones estatales centralizadas y coherentes en el siglo XVI, y en conjunción con desarrollos tecnológicos como la imprenta o el transporte marítimo de larga distancia, aparece en la Edad Moderna una creciente literatura que tiene por objetivo examinar el prbbtema del buen gobierno. Esto es, cómo asegurar la supervivencia de un estado, cuando otros estados aumentan su tamaño e influencia y a menudo, lo hacen por medio de la guerra. Estados de mayor tamaño, y con mayores amenazas por parte de otros estados rivales necesitaban imperiosamente políticas económicas que permitiesen financiar guerras y aumentar la población bajo su dominio. Surgen así variadas escuelas de pensamiento y autores que proponen al gobernante de tumo refonnas y actuaciones con contenido económico, en la línea de lo que hoy llamaríamos política económica. En la Edad Moderna y en Europa, destacan en este periodo obras de autores como Maquiavelo, los arbitristas españolas, la Escuela de Salamanca, los mercantilistas del siglo XVI (y sus críticos posteriores como PETTY o MANDEVILLE). En los años previos a la Revolución Industrial, los fisiócratas franceses y los cameralistas del norte de Europa desarrollan una amplia literatura que influye fuertemente en STEUART\footnote{JAMES STEUART publicó en 1767 la obra \textit{An Inquiry into the Principles of Political Economy}. Precede en nueve años al \textit{An Inquiry into the Nature and Causes of the Wealth of Nations} de ADAM SMITH. y ha sido referida a menudo como predecesora. Estudios históricos mostrarían que ADAM SMITH conocía esta obra de STEUART a pesar de no incluir referencia alguna en las suyas propias.} y ADAM SMITH. Y es con la obra de este último autor, \textit{Una investigación sobre la naturaleza y las causas de la riqueza de las naciones} con la que el enfoque del buen gobierno alcanza su culmen.

\textsc{Economía como descrubrimiento de leyes naturales}

La segunda mitad del siglo XVII y el siglo XVIII fueron testigos de un renovado impulso en lasciencias naturales. La física clásica de Newton, la invención del cálculo diferencial, el método deductivo de Descartes y avances generalizados en otras disciplinas como la biología, la astronomía o la química, influyeron fuertemente en las ciencias sociales. Así, si había sido posible explicar fenómenos como la gravedad, el movimiento de los cuerpos celestes o la circulación sanguínea a partir de leyes universales que describían la naturaleza de las cosas, también debería ser posible caracterizar el comportamiento humano en forma de leyes innatas de cumplimiento generalizado. En este contexto, la obra de los principales autores del periodo clásico tiene por objeto principal la búsqueda de esas leyes que expliquen el intercambio humano. 

Las obras de los principales autores del clasicismo británico tras Adam Smith están imbuidas de esta forma de entender el objeto de la ciencia económica. Así, la Teoría del Valor-Trabajo de Ricardo y Malthus, las diferentes teorías sobre la circulación monetaria\footnote{Esencialmente, la \textit{currency school} y la \textit{banking school}. \textit{Grosso modo}, la \textit{currency school} establecía una relación de causalidad desde la oferta monetaria hacia el nivel de precios y la economía real. Por su parte, la \textit{banking school} postulaba una causalidad en sentido opuesto: el volumen de actividad económica acaba por determinar la oferta monetaria.} del siglo XIX, la Teoría de la Población, los modelos de determinación de los salarios o el mecanismo de flujo-especie -entre otros-, son ejemplos de esta fonna de entender la ciencia económica: qué leyes o mecanismos rigen los intercambios humanos, y cómo caracterizarlas de forma general para entender lo que sucedió en el pasado y predecir lo que sucederá en el futuro. Aunque en cierta medida la consideración moral de los fenómenos económicos o el buen gobierno siguen estando presentes en el pensamiento económico, tales objetivos se ven en gran medida superados por esta nueva fonna de concebir el qué debe estudiar la economía.

\textsc{Economía como estudio de fenómenos psicológicos y morales que determinan producción}

A mediados del siglo XIX, en el periodo clásico tardío, JOHN STUART MILL propuso una matización no trivial de la concepción de la economía como la búsqueda de leyes naturales. Su forma de entender el objeto de la economía trata de incorporar la dimensión psicológica y política de la distribución y producción de los recursos. En la obra de este autor, el objeto de la economía no es alcanzar el buen gobierno de una nación (en términos del autor, esto sería confundir un arte con una ciencia), ni puede limitarse al descubrimiento de leyes más o menos generales. Esto es así porque la producción y la distribución de los recursos dependen del desarrollo tecnológico (que considera exógeno a la economía) y, sobre todo, de la acción política, psicológica y ética de los seres humanos. 

Así, en STUART MILL, el objeto de la economía es descubrir leyes, pero no tanto leyes de carácter tecnológico o natural, sino leyes psicológicas o morales que acaban determinando la producción y la distribución de los recursos. 

\textsc{Economía como gestión de asuntos de la vida corriente}

En la última década del siglo XIX, ALFRED MARSHALL publica \textit{Principles of Economics}. La obra fue pionera en cuanto a ser una de las primeras en utilizar el término economía y no economía política, más habitual hasta entonces. Entre otras muchas aportaciones de esta obra, además de esa nueva denominación, destaca una concepción del objeto de la economía diferente al de los clásicos. Para MARSHALL, la economía es la ciencia del comportamiento humano en sus quehaceres diarios, cuando éstos tienen por objeto alcanzar el máximo bienestar.\footnote{Literalmente: economics is a study of mankind in the ordinary business; it examines that part of individual and social action which is most closely connected with the attainment and with the use of the material requisites of wellbeing.}

Dado que las necesidades y los deseos humanos cambian con el tiempo, el objeto de la economía es también comprender la evolución de esas necesidades, así como de la propia evolución humana. 

\textsc{Economía como gestión de la escasez} 

La concepción de la economía de MARSHALL sentó las bases de la definición que LIONEL ROBBINS y PAUL SAMUELSON consolidarían posteriormente de manera casi canónica. 

Si el objeto de la economía es examinar las decisiones humanas en sus que haceres diarios, ¿cuáles son tales que haceres que merecen ser objeto de estudio por los economistas? LIONEL ROBBINS propuso una respuesta: 

La economía es el estudio del comportamiento humano cuando trata de gestionar recursos escasos con usos alternativos, para satisfacer necesidades humanas. 

PAUL SAMUELSON matizó después ligeramente esta definición, aunque manteniéndola en lo esencial. Y es ésta la definición que subyace a la microeconomía moderna: dada una renta, ¿cómo la distribuirá un consumidor entre diferentes bienes que puede adquirir?, ¿cómo repartirá su tiempo un trabajador-consumidor que decide cuanto tiempo dedicar al ocio y cuánto dedicarlo a trabajar para aumentar su consumo?, ¿cómo distribuirá temporalmente el estado la senda del gasto público, dados unos ingresos tributarios limitados? 

LIPSEY planteó posteriormente una matización a esta definición. Así, la economía no se limita a gestionar recursos escasos, si no que más bien trata de aprovechar de manera óptima los recursos disponibles, que pueden estar mal asignados. Si bien se trata de un matiz que no altera lá sustancia de la definición de ROBBINS y SAMUELSON, introduce un matiz que conecta con el objeto de los modelos macroeconómicos de corte keynesiano: si existe capacidad productiva sin aprovechar, ¿qué políticas económicas pueden llevar a la plena utilización de aquellos, que son escasos, pero en un momento determinado se encuentran ociosos?

\textsc{Economía como ciencia de la decisión óptima}

A partir de los años 60, el conjunto de métodos utilizados por los economistas para caracterizar la decisión humana pasa a aplicarse a un amplio rango de fenómenos humanos que trascienden las decisiones de producción, distribución y consumo. Así, la economía se convierte en una ciencia de la decisión óptima, en la que el objetivo es caracterizar el comportamiento tendente a alcanzar la satisfacción de unos objetivos definidos por unas preferencias definidas previamente. 

En esta línea, la economía se expande a otros dominios de las ciencias sociales: GARY BECKER caracteriza en términos microeconómicos decisiones tomadas en el seno de la familia tales como el matrimonio, el número de hijos o el comportamiento, así como la criminalidad, entre otros. ARROW caracteriza en términos matemáticos los límites a la formulación de constituciones políticas a partir de las preferencias individuales de los miembros de un grupo. BUCHANAN, DOWNS o TULLOCK aplican métodos similares al   proceso político y a la toma de decisiones en democracias representativas (y no sólo). Y numerosos otros autores adoptan este enfoque a problemas tan variados como la elección de carteras, la organización de empresas, las decisiones de marketing, el diseño de subastas o la política monetaria.

\anexo{Paradigma científico de KUHN en economía}\label{anx:paradigmaKUHN}
Aunque KUHN centró la atención en las ciencias naturales, influyó significativamente en las ciencias sociales, como la economía. En los análisis metodológicos de las décadas de 1970 y 1980 aparece repetidamente el término paradigma.

\textsc{Aplicación a los economistas marginalistas: la revolución marginalista}

Existe controversia en entender la obra de los economistas marginalistas como una revolución en términos de KUHN:
\begin{la}
  \item Por un lado, las implicaciones de política económica son análogas a las de sus predecesores de la escuela clásica (apuesta por el laissez-faire);
  \item Sin embargo, supone una importante revolución metodológica en dos aspectos:
  \begin{lnum}
    \item Individualismo metodológico: Los fenómenos sociales están explicados por comportamientos individuales. Esto contrasta con los clásicos, cuyos sujetos de análisis eran países o clases sociales.
    \item Empleo de herramientas de cálculo diferencial (formalización matemática a través de problemas de optimización): Esto contrasta con el enfoque de los clásicos, que hacían un mayor hincapié en el razonamiento verbal haciendo uso de obras fundamentalmente ensayísticas y una aproximación de economía política.
  \end{lnum}
\end{la}

\textsc{Aplicación a la Teoría Keyneisana}

Podríamos considerar que la aparición de la obra de KEYNES constituye en la ciencia económica un buen ejemplo de lo que KUHN denomina \textit{revolución científica}. La aparición del paradigma keynesiano ofrecía una respuesta a las anomalías que experimentaba el paradigma neoclásico y que se materializaban en las dificultades para ofrecer una explicación completa y convincente a los
efectos de la Gran Depresión.

No obstante, debemos señalar que la revolución científica que predecía KUHN no se produjo del modo que presagiaba dicho autor. En efecto, a pesar de las anomalías acumuladas por el paradigma neoclásico éste no quedó reemplazado completamente por el keynesiano sino que ambos coexistieron y llegaron a quedar integrados bajo el marco de la Síntesis Neoclásica.

\textsc{Aplicación de la Nueva Macroeconomía Clásica}

Igualmente, podemos entender la obra de los autores de la Nueva Macroeconomía Clásica como una revolución en términos de KUHN, iniciada por LUCAS y de algún modo redirigida y consolidada por KYDLAND y PRESCOTT, con cambios radicales en los temas de investigación y las herramientas metodológicas, dando lugar a distintas implicaciones de política económica.

En este sentido, cabe mencionar dos de las contribuciones de esta escuela por las que podemos considerarla una revolución en términos de KUHN: (i) la crítica de Lucas; y, (ii) y los modelos de Equilibrio General Dinámicos y Estocásticos (EGDE).


Crítica de Lucas (1976): LUCAS insiste en que no debe haber ninguna separación entre los principios metodológicos sobre los que se construye la teoría microeconómica y aquellos sobre los que se construye la teoría macroeconómica. Parte de que los modelos de la Síntesis Neoclásica hacían un buen trabajo predictivo, pero eran un fracaso total en la evaluación de políticas alternativas. El problema era el empleo de modelos macroeconométricos que no tenían en cuenta el comportamiento optimizador de los agentes y que se basan en puras relaciones de agregados. Para LUCAS: (i) los agentes toman decisiones optimizadoras; (ii) las reglas de decisión óptima cambian con cambios en las políticas económicas; (iii) cualquier cambio en las políticas va sistemáticamente a alterar la estructura de los modelos econométricos.

Los modelos de la Síntesis Neoclásica no tenían esto en cuenta y los modelos eran puras relaciones estadísticas que recogían el comportamientos pasados de los individuos. Esto provocaba que los coeficientes no tuvieran en cuenta posibles cambios en el período de políticas económicas. Es decir, los coeficientes son invariantes ante cambios en la orientación de las políticas económicas, cuando para LUCAS esto no es así porque los agentes reaccionan a cambios en las políticas económicas. En otras palabras, esos coeficientes estimados no son parámetros exógenos, sino que son variables endógenas sensibles a cambios en las políticas. Por tanto, había que especificar cuidadosamente la función objetivo, la restricción e hipótesis de comportamiento de los agentes para ver cómo los agentes se comportan (i.e. microfundamentar los modelos).\footnote{La microfundamentación es una condición sine qua non ya que sólo modelos profundamente estructurales que partan de los fundamentos de la economía son capaces de proporcionar una base sólida para la evaluación de políticas alternativas.}

\textsc{Modelos de Equilibrio General, Dinámicos y Estocásticos}

Se traga de modelos de agente representativo (una versión simplificada y tratable del modelo de Equlibrio General). 

Los agentes representativos en el modelo RBC básico son consumidores y empresas, con sendos problemas de optimización:
\begin{la}
  \item Los hogares maximizan su utilidad que depende del consumo (positivamente) y del trabajo (negativamente), se suele asumir separabilidad entre ambas por conveniencia, y están restringidos por una restricción de recursos, pues obtienen su renta de ofertar (a) trabajo a cambio de un salario real y (b) capital o ahorro a cambio de una rentabilidad real. Dicha renta se utiliza para adquirir bienes de consumo presente o para ahorrar/incrementar su riqueza; por otro lado,
  \item Las empresas operan en un entorno de competencia perfecta y cuya tecnología se resume en una función clave: la función de producción agregada, que expresa la cantidad total de producto (homogéneo) en función de las cantidades de los factores primarios (trabajo y capital) que demanda a los hogares.
\end{la}

Además, no suelen poseer sector público dado que no hay fallos de mercado de ningún tipo, por lo que los ciclos son respuestas óptimas de los agentes a los shocks e intervenir para combatirlos sería ineficiente (el equilibrio competitivo alcanzado es en todo momento óptimo de Pareto). La recomendación de política económica es desaconsejar las políticas macroeconómicas estabilizadoras.


Dinámico (D). Los agentes optimizan para varios períodos, por lo que maximizan utilidad/beneficios presentes y futuros descontados\footnote{No necesariamente forward looking. Hay modelos con retardos}\textsuperscript{,}No necesariamente forward looking. Hay modelos con retardos.
▫
Estocástico (E): Se produce un contexto de información imperfecta en que se incluyen perturbaciones exógenas. Esto obliga a aplicar la teoría de la utilidad esperada, de modo que los consumidores maximizan la esperanza de sus utilidades futuras, formando dichas expectativas racionalmente (hipótesis de expectativas racionales (HER)).


\anexo{Método deductivo y Método inductivo}\label{anx:espitemologia}
 

\end{document}